\documentclass[11pt,a4paper]{article}
\usepackage[utf8]{inputenc}
\usepackage[T1]{fontenc}
\usepackage[spanish,es-minimal,es-noshorthands]{babel}
\AtBeginDocument{\renewcommand{\sin}{\operatorname{sen}}\renewcommand{\arcsin}{\operatorname{arcsen}}}
\usepackage{lmodern}
\usepackage[sfdefault,lining]{FiraSans}
\usepackage[a4paper,top=2.3cm,bottom=2.2cm,left=1.8cm,right=1.8cm,headheight=1.3cm,headsep=0.4cm,footskip=1cm]{geometry}
\usepackage{amsmath,amssymb}
\usepackage{icomma}
\usepackage{graphicx}
\usepackage[table]{xcolor}
\usepackage{tikz}
\usepackage{fontawesome5}
\usepackage{ccicons}
\usepackage{fancyhdr}
\usepackage{lastpage}
\usepackage{tasks}
\usepackage{needspace}
\usepackage[most]{tcolorbox}
\usepackage{array}
\usepackage[hidelinks]{hyperref}
\AtBeginDocument{\hypersetup{pdftitle={Soluciones de la ficha de la unidad 6: Funciones. Características. 4º ESO},pdfsubject={Matemáticas A 4º ESO},pdfkeywords={funciones, variables, imagen, antiimagen, dominio, recorrido, puntos de corte, continuidad, discontinuidades, crecimiento, máximos y mínimos, periodicidad, interpretación de gráficos, construcción de funciones}}}

\definecolor{navy}{HTML}{1B2748}
\definecolor{accent}{HTML}{E8603A}
\definecolor{teal}{HTML}{1F9E90}
\definecolor{soft}{HTML}{EEF1F8}
\definecolor{muted}{HTML}{5D6685}

\setlength{\parindent}{0pt}
\setlength{\parskip}{3pt}
\settasks{label=\alph*), label-format=\bfseries\color{navy}, label-width=1.3em, item-indent=1.8em, column-sep=0.6em, after-item-skip=5pt, before-skip=3pt, after-skip=4pt}

% ---------- capçalera i peu ----------
\newcommand{\llicencia}{\href{https://creativecommons.org/licenses/by-nc-sa/4.0/deed.es}{\ccbyncsa}\ \ Licencia Creative Commons Atribución-NoComercial-CompartirIgual 4.0 Internacional (CC BY-NC-SA 4.0)}
\pagestyle{fancy}
\fancyhf{}
\fancyhead[L]{\raisebox{-0.35\height}{\includegraphics[height=0.95cm]{logo-gva.png}}}
\fancyhead[C]{\small\color{muted}Soluciones de la ficha · Unidad 6: Funciones. Características · 4\textsuperscript{o} ESO A}
\fancyhead[R]{\raisebox{-0.35\height}{\includegraphics[height=0.95cm]{logo-vermellar.png}}}
\fancyfoot[L]{\scriptsize\color{muted}\llicencia}
\fancyfoot[R]{\small\color{muted}\thepage}
\renewcommand{\headrulewidth}{0.4pt}
\renewcommand{\headrule}{\color{muted!40}\hrule height \headrulewidth}
\fancypagestyle{primera}{\fancyhf{}\renewcommand{\headrulewidth}{0pt}
  \fancyfoot[L]{\scriptsize\color{muted}\llicencia}
  \fancyfoot[R]{\small\color{muted}\thepage}}

% ---------- nivells i exercicis ----------
\newcommand{\estrelles}[1]{{\color{accent}\ifcase#1\or\faStar\faStar[regular]\faStar[regular]\or\faStar\faStar\faStar[regular]\or\faStar\faStar\faStar\fi}}
\newcounter{exer}
\newcommand{\ex}[1]{\par\needspace{4\baselineskip}\medskip\refstepcounter{exer}\noindent{\color{navy}\bfseries\theexer.}\ \ {\small\estrelles{#1}}\ \ }

\newcommand{\seccio}[2]{\par\needspace{16\baselineskip}\bigskip
  \begin{tcolorbox}[colback=navy,colframe=navy,coltext=white,arc=2mm,boxrule=0pt,left=3mm,right=3mm,top=1.2mm,bottom=1.2mm]
  {\large\bfseries #1\quad #2}
  \end{tcolorbox}\vspace{-2pt}}

\newcommand{\exf}[1]{\par\medskip\refstepcounter{exer}\noindent{\color{navy}\bfseries\theexer.}\ \ {\small\estrelles{#1}}\ \ }
\newcolumntype{C}[1]{>{\centering\arraybackslash}p{#1}}
% quadre de recordatori (només a la fitxa d'exercicis)
\newcommand{\recorda}[1]{\par\smallskip\begin{tcolorbox}[colback=accent!7,colframe=accent!55,arc=2mm,boxrule=0.5pt,left=2.5mm,right=2.5mm,top=1mm,bottom=1mm]\small{\color{accent}\faLightbulb[regular]}\ \textbf{Recuerda:}\ #1\end{tcolorbox}}
\usetikzlibrary{calc}
\usepackage{enumitem}
\definecolor{violet}{HTML}{7C3AED}
\definecolor{orange}{HTML}{EA8A1F}
\newcommand{\res}[1]{\boxed{#1}}\newcommand{\rest}[1]{\fbox{#1}}

\def\fmneed{11}% que el títol d'una secció no quede sol al final del full
% =====================================================================
% Capçalera FentMates de les fitxes: un color per a cada curs.
% Els colors són els de la pàgina inicial (LEVELS a assets/js/materials.js);
% el logo ve de _eines/logo/pdf/ (python3 _eines/logo/genera.py).
%
% Ús, al final del preàmbul de la unitat:
%   % =====================================================================
% Capçalera FentMates de les fitxes: un color per a cada curs.
% Els colors són els de la pàgina inicial (LEVELS a assets/js/materials.js);
% el logo ve de _eines/logo/pdf/ (python3 _eines/logo/genera.py).
%
% Ús, al final del preàmbul de la unitat:
%   % =====================================================================
% Capçalera FentMates de les fitxes: un color per a cada curs.
% Els colors són els de la pàgina inicial (LEVELS a assets/js/materials.js);
% el logo ve de _eines/logo/pdf/ (python3 _eines/logo/genera.py).
%
% Ús, al final del preàmbul de la unitat:
%   \input{../comu/fentmates.tex}
%   \fmcurs{1eso}
% i, al començament del document:
%   \capcalera{Matemàtiques}{1r ESO}{Fitxa d'exercicis · Unitat 2}{Divisibilitat}
% Opcional abans de \input: \def\fmneed{14} (espai mínim abans de cada secció).
% =====================================================================
\makeatletter
\def\fm@logodir{../../logo/pdf}
\providecommand{\fmneed}{10}
\def\fm@cursnom{}

% clau del curs, color i color del text damunt del color
\newcommand{\fm@curs}[3]{\@namedef{fm@color@#1}{#2}\@namedef{fm@text@#1}{#3}}
\fm@curs{1eso}{E6B84A}{navy}
\fm@curs{2eso}{3B82F6}{white}
\fm@curs{3eso}{FF7A50}{white}
\fm@curs{4eso-a}{A78BFA}{white}
\fm@curs{4eso-b}{6D28D9}{white}
\fm@curs{1batx-mi}{9B1C1C}{white}
\fm@curs{1batx-ccss}{EF6B6B}{white}
\fm@curs{2batx-mii}{1E6B47}{white}
\fm@curs{2batx-ccss}{2BB3A3}{white}

\newcommand{\fmcurs}[1]{%
  \@ifundefined{fm@color@#1}{\PackageError{fentmates}{Curs desconegut: #1}{}}{}%
  \edef\fm@tmp{\noexpand\definecolor{curs}{HTML}{\@nameuse{fm@color@#1}}%
               \noexpand\colorlet{curstext}{\@nameuse{fm@text@#1}}}\fm@tmp
  \colorlet{cursfosc}{curs!72!black}%
  \def\fm@logo{\fm@logodir/vertical-#1.pdf}%
  \def\fm@marca{\fm@logodir/marca-#1.pdf}}

% ---------- capçalera de la primera pàgina ----------
% #1 matèria (petit), #2 curs, #3 tipus de document i unitat, #4 títol
\newcommand{\capcalera}[4]{%
  \gdef\fm@cursnom{#2}%
  \par\noindent
  \begin{tcolorbox}[colback=curs,colframe=curs,boxrule=0pt,arc=2.2mm,width=3.5cm,height=2.15cm,
      valign=center,left=3mm,right=2mm,top=1mm,bottom=1mm,nobeforeafter,box align=center]
    \color{curstext}{\footnotesize #1}\par\vspace{1pt}{\fontsize{19}{21}\selectfont\bfseries #2\par}%
  \end{tcolorbox}\hspace{5mm}%
  \begin{minipage}[c]{\dimexpr\linewidth-3.5cm-5mm-2.9cm\relax}
    {\color{cursfosc}\small\bfseries\MakeUppercase{#3}}\par\vspace{2pt}
    {\color{navy}\fontsize{20}{23}\selectfont\bfseries #4\par}
  \end{minipage}\hfill
  \begin{minipage}[c]{2.6cm}\raggedleft\includegraphics[width=2.6cm]{\fm@logo}\end{minipage}\par
  \vspace{2pt}}

% ---------- barres de secció del color del curs ----------
\renewcommand{\seccio}[2]{\par\needspace{\fmneed\baselineskip}\bigskip
  \begin{tcolorbox}[colback=curs,colframe=curs,coltext=curstext,arc=2mm,boxrule=0pt,left=3mm,right=3mm,top=1.2mm,bottom=1.2mm]
  {\large\bfseries #1\quad #2}
  \end{tcolorbox}\vspace{-2pt}}

% ---------- estrelles de dificultat del color del curs (a Batxillerat no se'n posen) ----------
\@ifundefined{estrelles}{}{%
  \renewcommand{\estrelles}[1]{{\color{curs}\ifcase#1\or\faStar\faStar[regular]\faStar[regular]\or\faStar\faStar\faStar[regular]\or\faStar\faStar\faStar\fi}}}

% ---------- capçalera i peu de les pàgines ----------
% primera pàgina: logos institucionals al peu
\fancypagestyle{primera}{\fancyhf{}\renewcommand{\headrulewidth}{0pt}
  \fancyfoot[L]{\raisebox{-0.32\height}{\includegraphics[height=0.78cm]{logo-gva.png}}\hspace{5mm}%
    \raisebox{-0.32\height}{\includegraphics[height=0.78cm]{logo-vermellar.png}}}
  \fancyfoot[R]{\scriptsize\color{muted}\href{https://creativecommons.org/licenses/by-nc-sa/4.0/deed.ca}{\ccbyncsa}\ \ CC BY-NC-SA 4.0 · fentmates.org\qquad\small\thepage}}
% la resta: el símbol a l'esquerra i el curs a la dreta, amb el títol al centre com abans
\fancyhead[L]{\raisebox{-0.3\height}{\includegraphics[height=0.62cm]{\fm@marca}}}
\fancyhead[R]{\small\bfseries\color{cursfosc}\fm@cursnom}
\renewcommand{\headrule}{\color{curs!45}\hrule height \headrulewidth}
\makeatother

%   \fmcurs{1eso}
% i, al començament del document:
%   \capcalera{Matemàtiques}{1r ESO}{Fitxa d'exercicis · Unitat 2}{Divisibilitat}
% Opcional abans de \input: \def\fmneed{14} (espai mínim abans de cada secció).
% =====================================================================
\makeatletter
\def\fm@logodir{../../logo/pdf}
\providecommand{\fmneed}{10}
\def\fm@cursnom{}

% clau del curs, color i color del text damunt del color
\newcommand{\fm@curs}[3]{\@namedef{fm@color@#1}{#2}\@namedef{fm@text@#1}{#3}}
\fm@curs{1eso}{E6B84A}{navy}
\fm@curs{2eso}{3B82F6}{white}
\fm@curs{3eso}{FF7A50}{white}
\fm@curs{4eso-a}{A78BFA}{white}
\fm@curs{4eso-b}{6D28D9}{white}
\fm@curs{1batx-mi}{9B1C1C}{white}
\fm@curs{1batx-ccss}{EF6B6B}{white}
\fm@curs{2batx-mii}{1E6B47}{white}
\fm@curs{2batx-ccss}{2BB3A3}{white}

\newcommand{\fmcurs}[1]{%
  \@ifundefined{fm@color@#1}{\PackageError{fentmates}{Curs desconegut: #1}{}}{}%
  \edef\fm@tmp{\noexpand\definecolor{curs}{HTML}{\@nameuse{fm@color@#1}}%
               \noexpand\colorlet{curstext}{\@nameuse{fm@text@#1}}}\fm@tmp
  \colorlet{cursfosc}{curs!72!black}%
  \def\fm@logo{\fm@logodir/vertical-#1.pdf}%
  \def\fm@marca{\fm@logodir/marca-#1.pdf}}

% ---------- capçalera de la primera pàgina ----------
% #1 matèria (petit), #2 curs, #3 tipus de document i unitat, #4 títol
\newcommand{\capcalera}[4]{%
  \gdef\fm@cursnom{#2}%
  \par\noindent
  \begin{tcolorbox}[colback=curs,colframe=curs,boxrule=0pt,arc=2.2mm,width=3.5cm,height=2.15cm,
      valign=center,left=3mm,right=2mm,top=1mm,bottom=1mm,nobeforeafter,box align=center]
    \color{curstext}{\footnotesize #1}\par\vspace{1pt}{\fontsize{19}{21}\selectfont\bfseries #2\par}%
  \end{tcolorbox}\hspace{5mm}%
  \begin{minipage}[c]{\dimexpr\linewidth-3.5cm-5mm-2.9cm\relax}
    {\color{cursfosc}\small\bfseries\MakeUppercase{#3}}\par\vspace{2pt}
    {\color{navy}\fontsize{20}{23}\selectfont\bfseries #4\par}
  \end{minipage}\hfill
  \begin{minipage}[c]{2.6cm}\raggedleft\includegraphics[width=2.6cm]{\fm@logo}\end{minipage}\par
  \vspace{2pt}}

% ---------- barres de secció del color del curs ----------
\renewcommand{\seccio}[2]{\par\needspace{\fmneed\baselineskip}\bigskip
  \begin{tcolorbox}[colback=curs,colframe=curs,coltext=curstext,arc=2mm,boxrule=0pt,left=3mm,right=3mm,top=1.2mm,bottom=1.2mm]
  {\large\bfseries #1\quad #2}
  \end{tcolorbox}\vspace{-2pt}}

% ---------- estrelles de dificultat del color del curs (a Batxillerat no se'n posen) ----------
\@ifundefined{estrelles}{}{%
  \renewcommand{\estrelles}[1]{{\color{curs}\ifcase#1\or\faStar\faStar[regular]\faStar[regular]\or\faStar\faStar\faStar[regular]\or\faStar\faStar\faStar\fi}}}

% ---------- capçalera i peu de les pàgines ----------
% primera pàgina: logos institucionals al peu
\fancypagestyle{primera}{\fancyhf{}\renewcommand{\headrulewidth}{0pt}
  \fancyfoot[L]{\raisebox{-0.32\height}{\includegraphics[height=0.78cm]{logo-gva.png}}\hspace{5mm}%
    \raisebox{-0.32\height}{\includegraphics[height=0.78cm]{logo-vermellar.png}}}
  \fancyfoot[R]{\scriptsize\color{muted}\href{https://creativecommons.org/licenses/by-nc-sa/4.0/deed.ca}{\ccbyncsa}\ \ CC BY-NC-SA 4.0 · fentmates.org\qquad\small\thepage}}
% la resta: el símbol a l'esquerra i el curs a la dreta, amb el títol al centre com abans
\fancyhead[L]{\raisebox{-0.3\height}{\includegraphics[height=0.62cm]{\fm@marca}}}
\fancyhead[R]{\small\bfseries\color{cursfosc}\fm@cursnom}
\renewcommand{\headrule}{\color{curs!45}\hrule height \headrulewidth}
\makeatother

%   \fmcurs{1eso}
% i, al començament del document:
%   \capcalera{Matemàtiques}{1r ESO}{Fitxa d'exercicis · Unitat 2}{Divisibilitat}
% Opcional abans de \input: \def\fmneed{14} (espai mínim abans de cada secció).
% =====================================================================
\makeatletter
\def\fm@logodir{../../logo/pdf}
\providecommand{\fmneed}{10}
\def\fm@cursnom{}

% clau del curs, color i color del text damunt del color
\newcommand{\fm@curs}[3]{\@namedef{fm@color@#1}{#2}\@namedef{fm@text@#1}{#3}}
\fm@curs{1eso}{E6B84A}{navy}
\fm@curs{2eso}{3B82F6}{white}
\fm@curs{3eso}{FF7A50}{white}
\fm@curs{4eso-a}{A78BFA}{white}
\fm@curs{4eso-b}{6D28D9}{white}
\fm@curs{1batx-mi}{9B1C1C}{white}
\fm@curs{1batx-ccss}{EF6B6B}{white}
\fm@curs{2batx-mii}{1E6B47}{white}
\fm@curs{2batx-ccss}{2BB3A3}{white}

\newcommand{\fmcurs}[1]{%
  \@ifundefined{fm@color@#1}{\PackageError{fentmates}{Curs desconegut: #1}{}}{}%
  \edef\fm@tmp{\noexpand\definecolor{curs}{HTML}{\@nameuse{fm@color@#1}}%
               \noexpand\colorlet{curstext}{\@nameuse{fm@text@#1}}}\fm@tmp
  \colorlet{cursfosc}{curs!72!black}%
  \def\fm@logo{\fm@logodir/vertical-#1.pdf}%
  \def\fm@marca{\fm@logodir/marca-#1.pdf}}

% ---------- capçalera de la primera pàgina ----------
% #1 matèria (petit), #2 curs, #3 tipus de document i unitat, #4 títol
\newcommand{\capcalera}[4]{%
  \gdef\fm@cursnom{#2}%
  \par\noindent
  \begin{tcolorbox}[colback=curs,colframe=curs,boxrule=0pt,arc=2.2mm,width=3.5cm,height=2.15cm,
      valign=center,left=3mm,right=2mm,top=1mm,bottom=1mm,nobeforeafter,box align=center]
    \color{curstext}{\footnotesize #1}\par\vspace{1pt}{\fontsize{19}{21}\selectfont\bfseries #2\par}%
  \end{tcolorbox}\hspace{5mm}%
  \begin{minipage}[c]{\dimexpr\linewidth-3.5cm-5mm-2.9cm\relax}
    {\color{cursfosc}\small\bfseries\MakeUppercase{#3}}\par\vspace{2pt}
    {\color{navy}\fontsize{20}{23}\selectfont\bfseries #4\par}
  \end{minipage}\hfill
  \begin{minipage}[c]{2.6cm}\raggedleft\includegraphics[width=2.6cm]{\fm@logo}\end{minipage}\par
  \vspace{2pt}}

% ---------- barres de secció del color del curs ----------
\renewcommand{\seccio}[2]{\par\needspace{\fmneed\baselineskip}\bigskip
  \begin{tcolorbox}[colback=curs,colframe=curs,coltext=curstext,arc=2mm,boxrule=0pt,left=3mm,right=3mm,top=1.2mm,bottom=1.2mm]
  {\large\bfseries #1\quad #2}
  \end{tcolorbox}\vspace{-2pt}}

% ---------- estrelles de dificultat del color del curs (a Batxillerat no se'n posen) ----------
\@ifundefined{estrelles}{}{%
  \renewcommand{\estrelles}[1]{{\color{curs}\ifcase#1\or\faStar\faStar[regular]\faStar[regular]\or\faStar\faStar\faStar[regular]\or\faStar\faStar\faStar\fi}}}

% ---------- capçalera i peu de les pàgines ----------
% primera pàgina: logos institucionals al peu
\fancypagestyle{primera}{\fancyhf{}\renewcommand{\headrulewidth}{0pt}
  \fancyfoot[L]{\raisebox{-0.32\height}{\includegraphics[height=0.78cm]{logo-gva.png}}\hspace{5mm}%
    \raisebox{-0.32\height}{\includegraphics[height=0.78cm]{logo-vermellar.png}}}
  \fancyfoot[R]{\scriptsize\color{muted}\href{https://creativecommons.org/licenses/by-nc-sa/4.0/deed.ca}{\ccbyncsa}\ \ CC BY-NC-SA 4.0 · fentmates.org\qquad\small\thepage}}
% la resta: el símbol a l'esquerra i el curs a la dreta, amb el títol al centre com abans
\fancyhead[L]{\raisebox{-0.3\height}{\includegraphics[height=0.62cm]{\fm@marca}}}
\fancyhead[R]{\small\bfseries\color{cursfosc}\fm@cursnom}
\renewcommand{\headrule}{\color{curs!45}\hrule height \headrulewidth}
\makeatother

\fmcurs{4eso-a}

\begin{document}
\thispagestyle{primera}

\capcalera{Matemáticas A}{4\textsuperscript{o} ESO}{Soluciones de la ficha · Unidad 6}{\raggedright Funciones. Características}
\vspace{8pt}

\begin{tcolorbox}[colback=soft,colframe=soft,arc=2mm,left=3mm,right=3mm,top=1.5mm,bottom=1.5mm]
\small Intenta resolver cada ejercicio antes de mirar su solución. En los ejercicios de «dibuja una función…» hay muchas respuestas correctas: comprueba que la tuya cumpla todas las condiciones. Puedes comprobar dominios y cortes con las herramientas interactivas de los apuntes de la unidad.
\end{tcolorbox}

% =====================================================
\seccio{1}{Concepto de función. Imágenes y antiimágenes}

\exf{1} \begin{tasks}(4)
  \task $\res{\text{Sí}}$
  \task $\res{\text{No}}$
  \task $\res{\text{Sí}}$
  \task $\res{\text{No}}$
\end{tasks}Las gráficas b) y d) tienen valores de $x$ con dos imágenes: hay rectas verticales que las cortan dos veces.

\ex{1} \begin{tasks}(1)
  \task $\res{\text{Sí}}$
  \task $\res{\text{No}}$: el 4 tiene dos raíces, $2$ y $-2$.
  \task $\res{\text{Sí}}$
  \task $\res{\text{No}}$: un alumno puede tener dos o más hermanos (o ninguno).
\end{tasks}

\ex{1} \begin{tasks}(1)
  \task Independiente: \rest{número de libretas}; dependiente: \rest{precio}. \rest{No se unen}: no se pueden comprar 2,5 libretas.
  \task Independiente: \rest{tiempo}; dependiente: \rest{altura}. \rest{Se unen}: la planta crece continuamente.
  \task Independiente: \rest{velocidad}; dependiente: \rest{tiempo}. \rest{Se unen}: la velocidad puede tomar cualquier valor.
\end{tasks}

\ex{2} \begin{tasks}(1)
  \task $f(-2)=4+4-3=\res{5}$; \ $f(0)=\res{-3}$; \ $f\left(\frac12\right)=\frac14-1-3=\res{-\frac{15}{4}}$
  \task $x^2-2x-3=0\Rightarrow\res{x=-1,\ x=3}$; \ $x^2-2x+1=0\Rightarrow\res{x=1}$; \ $x^2-2x+3=0$, $\Delta=-8<0$: $-6$ \rest{no tiene antiimágenes}.
\end{tasks}

\ex{1} \begin{center}\renewcommand{\arraystretch}{1.25}\begin{tabular}{|c|C{1cm}|C{1cm}|C{1cm}|C{1cm}|C{1cm}|}\hline $x$ & $-2$ & $-1$ & $0$ & $1$ & $2$\\\hline $f(x)$ & $\res{-8}$ & $\res{-5}$ & $\res{-2}$ & $\res{1}$ & $\res{4}$\\\hline\end{tabular}\end{center}La gráfica es una recta. \ $3x-2=10\Rightarrow\res{x=4}$

\begin{center}\begin{tikzpicture}[x=0.467cm,y=0.23cm,baseline={(current bounding box.center)},font=\scriptsize]
\draw[muted!25,line width=.3pt] (-3,-9)--(-3,5);
\draw[muted!25,line width=.3pt] (-2,-9)--(-2,5);
\draw[muted!25,line width=.3pt] (-1,-9)--(-1,5);
\draw[muted!25,line width=.3pt] (0,-9)--(0,5);
\draw[muted!25,line width=.3pt] (1,-9)--(1,5);
\draw[muted!25,line width=.3pt] (2,-9)--(2,5);
\draw[muted!25,line width=.3pt] (3,-9)--(3,5);
\draw[muted!25,line width=.3pt] (-3,-9)--(3,-9);
\draw[muted!25,line width=.3pt] (-3,-8)--(3,-8);
\draw[muted!25,line width=.3pt] (-3,-7)--(3,-7);
\draw[muted!25,line width=.3pt] (-3,-6)--(3,-6);
\draw[muted!25,line width=.3pt] (-3,-5)--(3,-5);
\draw[muted!25,line width=.3pt] (-3,-4)--(3,-4);
\draw[muted!25,line width=.3pt] (-3,-3)--(3,-3);
\draw[muted!25,line width=.3pt] (-3,-2)--(3,-2);
\draw[muted!25,line width=.3pt] (-3,-1)--(3,-1);
\draw[muted!25,line width=.3pt] (-3,0)--(3,0);
\draw[muted!25,line width=.3pt] (-3,1)--(3,1);
\draw[muted!25,line width=.3pt] (-3,2)--(3,2);
\draw[muted!25,line width=.3pt] (-3,3)--(3,3);
\draw[muted!25,line width=.3pt] (-3,4)--(3,4);
\draw[muted!25,line width=.3pt] (-3,5)--(3,5);
\draw[-latex,muted] (-3,0)--(3,0);
\draw[-latex,muted] (0,-9)--(0,5);
\node[below=-1pt,text=muted] at (-3,0) {$-3$};
\node[below=-1pt,text=muted] at (-2,0) {$-2$};
\node[below=-1pt,text=muted] at (-1,0) {$-1$};
\node[below=-1pt,text=muted] at (1,0) {$1$};
\node[below=-1pt,text=muted] at (2,0) {$2$};
\node[below=-1pt,text=muted] at (3,0) {$3$};
\node[left=-1pt,text=muted] at (0,-8) {$-8$};
\node[left=-1pt,text=muted] at (0,-6) {$-6$};
\node[left=-1pt,text=muted] at (0,-4) {$-4$};
\node[left=-1pt,text=muted] at (0,-2) {$-2$};
\node[left=-1pt,text=muted] at (0,2) {$2$};
\node[left=-1pt,text=muted] at (0,4) {$4$};
\begin{scope}\clip (-3,-9) rectangle (3,5);
\draw[accent,line width=1.1pt,line join=round] (-2,-8)--(-1,-5);
\draw[accent,line width=1.1pt,line join=round] (-1,-5)--(0,-2);
\draw[accent,line width=1.1pt,line join=round] (0,-2)--(1,1);
\draw[accent,line width=1.1pt,line join=round] (1,1)--(2,4);
\filldraw[fill=accent,draw=accent,line width=.8pt] (-2,-8) circle (1.5pt);
\filldraw[fill=accent,draw=accent,line width=.8pt] (-1,-5) circle (1.5pt);
\filldraw[fill=accent,draw=accent,line width=.8pt] (0,-2) circle (1.5pt);
\filldraw[fill=accent,draw=accent,line width=.8pt] (1,1) circle (1.5pt);
\filldraw[fill=accent,draw=accent,line width=.8pt] (2,4) circle (1.5pt);
\end{scope}\end{tikzpicture}\end{center}


\ex{2} \begin{tasks}(4)
  \task $\res{f(x)=3x+5}$
  \task $\res{A(x)=x^2}$
  \task $\res{P(x)=1{,}2x}$
  \task $\res{f(x)=\dfrac{x^2}{2}}$
\end{tasks}

\ex{2} \begin{tasks}(1)
  \task $g(0)=\res{-\frac12}$, \ $g(3)=\res{4}$, \ $g(-1)=\res{0}$. \ $g(2)=\frac30$ \rest{no existe}: 2 no es del dominio.
  \task $\frac{x+1}{x-2}=2\Rightarrow x+1=2x-4\Rightarrow\res{x=5}$. \ $\frac{x+1}{x-2}=1\Rightarrow x+1=x-2\Rightarrow 1=-2$: el 1 \rest{no tiene antiimágenes}.
\end{tasks}


% =====================================================
\seccio{2}{Dominio y recorrido}

\exf{1} \begin{tasks}(2)
  \task $D=\res{[-1,\,3]}$, $R=\res{[-2,\,2]}$
  \task $D=\res{\mathbb{R}-\{-1\}}$, $R=\res{\mathbb{R}-\{2\}}$
  \task $D=\res{[-2,\,+\infty)}$, $R=\res{[-1,\,+\infty)}$
  \task $D=\res{[-4,\,2]}$, $R=\res{[-1,\,2)\cup\{3\}}$
\end{tasks}

\ex{1} \begin{tasks}(3)
  \task $\res{\mathbb{R}}$
  \task $\res{\mathbb{R}-\{5\}}$
  \task $\res{\mathbb{R}-\{-3\}}$
  \task $\res{\mathbb{R}-\{-3,\,3\}}$
  \task $\res{\mathbb{R}}$
  \task $\res{\mathbb{R}-\{0,\,6\}}$
\end{tasks}

\ex{2} \begin{tasks}(3)
  \task $\res{[-7,\,+\infty)}$
  \task $\res{[3,\,+\infty)}$
  \task $\res{(-\infty,\,2]}$
  \task $\res{(2,\,+\infty)}$
  \task $\res{(-\infty,\,4)}$
  \task $\res{[0,\,9)\cup(9,\,+\infty)}$
\end{tasks}

\ex{1} $f(1)=\sqrt{-2}$ no existe: $\res{1\notin D(f)}$. \ $f(3)=\sqrt0=0$: $\res{3\in D(f)}$. \ $f(7)=\sqrt4=2$: $\res{7\in D(f)}$. \ $x-3\geq0\Rightarrow D(f)=\res{[3,\,+\infty)}$

\ex{2} \begin{tasks}(2)
  \task una recta no horizontal toma todos los valores: $R=\res{\mathbb{R}}$
  \task $x^2\geq0\Rightarrow R=\res{[2,\,+\infty)}$
  \task $\sqrt{x-1}\geq0\Rightarrow R=\res{[3,\,+\infty)}$
  \task $-x^2\le0\Rightarrow R=\res{(-\infty,\,5]}$
\end{tasks}

\ex{3} Por ejemplo (hay muchas): \begin{tasks}(3)
  \task $\res{f(x)=\dfrac{1}{x-4}}$
  \task $\res{f(x)=\sqrt{x-2}}$
  \task $\res{f(x)=\dfrac{1}{\sqrt{5-x}}}$
\end{tasks}


% =====================================================
\seccio{3}{Puntos de corte con los ejes}

\exf{1} \begin{tasks}(2)
  \task $OY$: $\res{(0,\,-6)}$; $OX$: $\res{(3,\,0)}$
  \task $OY$: $\res{(0,\,12)}$; $OX$: $\res{(4,\,0)}$
  \task $OY$: $\res{(0,\,-16)}$; $OX$: $\res{(-4,\,0),\ (4,\,0)}$
  \task $OY$: $\res{(0,\,0)}$; $OX$: $\res{(-5,\,0),\ (0,\,0)}$
  \task $OY$: $\res{(0,\,4)}$; $OX$: $\res{(1,\,0),\ (4,\,0)}$
  \task $OY$: $\res{(0,\,5)}$; $OX$: no tiene
  \task $OY$: $\res{(0,\,-2)}$; $OX$: $\res{(-4,\,0)}$
  \task $OY$: $\res{(0,\,1)}$; $OX$: $\res{(-5,\,0)}$
\end{tasks}

\ex{1} $OX$: $\res{(-4,\,0),\ (-1,\,0),\ (2,\,0)}$; \ $OY$: $\res{(0,\,-1)}$.

\ex{2} $f(2)=4+2a-6=0\Rightarrow\res{a=1}$. \ $x^2+x-6=0\Rightarrow$ el otro corte es $\res{(-3,\,0)}$. \ $OY$: $\res{(0,\,-6)}$.

\ex{3} Si se anula en $-1$ y en $4$, es $f(x)=a(x+1)(x-4)$. \ $f(0)=-4a=-8\Rightarrow a=2$: \ $\res{f(x)=2(x+1)(x-4)=2x^2-6x-8}$

\ex{2} \begin{tasks}(1)
  \task $x(x^2-4)=x(x-2)(x+2)=0$: $OX$: $\res{(-2,\,0),\ (0,\,0),\ (2,\,0)}$; $OY$: $\res{(0,\,0)}$
  \task $x(x^2-x-6)=x(x-3)(x+2)=0$: $OX$: $\res{(-2,\,0),\ (0,\,0),\ (3,\,0)}$; $OY$: $\res{(0,\,0)}$
\end{tasks}


% =====================================================
\seccio{4}{Continuidad y tipos de discontinuidades}

\exf{1} \begin{tasks}(1)
  \task Discontinuidad $\res{\text{evitable}}$: los dos lados llegan a 3, pero $f(2)$ no existe. Sería continua si $\res{f(2)=3}$.
  \task Discontinuidad de $\res{\text{salto finito}}$: por la izquierda llega a 1 (y $f(2)=1$) y por la derecha a 3.
  \task $\res{\text{Continua}}$: los dos lados llegan a 1 y $f(2)=1$. Tiene un pico, pero se dibuja sin levantar el lápiz.
  \task Discontinuidad $\res{\text{asintótica}}$: por la izquierda se va a $-\infty$ y por la derecha a $+\infty$ (asíntota vertical $x=2$).
\end{tasks}

\ex{2} $x=-3$: $\res{\text{evitable}}$ (el punto está fuera de lugar). \ $x=0$: $\res{\text{salto finito}}$ (de 3 a 1). \ $x=2$: $\res{\text{salto infinito}}$ (por la izquierda llega a 0 y por la derecha se va a $+\infty$).

\ex{2} Por ejemplo (hay muchas):

\begin{center}\begin{tikzpicture}[x=0.46cm,y=0.46cm,baseline={(current bounding box.center)},font=\scriptsize]
\draw[muted!25,line width=.3pt] (-3,-2)--(-3,6);
\draw[muted!25,line width=.3pt] (-2,-2)--(-2,6);
\draw[muted!25,line width=.3pt] (-1,-2)--(-1,6);
\draw[muted!25,line width=.3pt] (0,-2)--(0,6);
\draw[muted!25,line width=.3pt] (1,-2)--(1,6);
\draw[muted!25,line width=.3pt] (2,-2)--(2,6);
\draw[muted!25,line width=.3pt] (3,-2)--(3,6);
\draw[muted!25,line width=.3pt] (4,-2)--(4,6);
\draw[muted!25,line width=.3pt] (5,-2)--(5,6);
\draw[muted!25,line width=.3pt] (6,-2)--(6,6);
\draw[muted!25,line width=.3pt] (7,-2)--(7,6);
\draw[muted!25,line width=.3pt] (-3,-2)--(7,-2);
\draw[muted!25,line width=.3pt] (-3,-1)--(7,-1);
\draw[muted!25,line width=.3pt] (-3,0)--(7,0);
\draw[muted!25,line width=.3pt] (-3,1)--(7,1);
\draw[muted!25,line width=.3pt] (-3,2)--(7,2);
\draw[muted!25,line width=.3pt] (-3,3)--(7,3);
\draw[muted!25,line width=.3pt] (-3,4)--(7,4);
\draw[muted!25,line width=.3pt] (-3,5)--(7,5);
\draw[muted!25,line width=.3pt] (-3,6)--(7,6);
\draw[-latex,muted] (-3,0)--(7,0);
\draw[-latex,muted] (0,-2)--(0,6);
\node[below=-1pt,text=muted] at (-3,0) {$-3$};
\node[below=-1pt,text=muted] at (-2,0) {$-2$};
\node[below=-1pt,text=muted] at (-1,0) {$-1$};
\node[below=-1pt,text=muted] at (1,0) {$1$};
\node[below=-1pt,text=muted] at (2,0) {$2$};
\node[below=-1pt,text=muted] at (3,0) {$3$};
\node[below=-1pt,text=muted] at (4,0) {$4$};
\node[below=-1pt,text=muted] at (5,0) {$5$};
\node[below=-1pt,text=muted] at (6,0) {$6$};
\node[below=-1pt,text=muted] at (7,0) {$7$};
\node[left=-1pt,text=muted] at (0,-2) {$-2$};
\node[left=-1pt,text=muted] at (0,-1) {$-1$};
\node[left=-1pt,text=muted] at (0,1) {$1$};
\node[left=-1pt,text=muted] at (0,2) {$2$};
\node[left=-1pt,text=muted] at (0,3) {$3$};
\node[left=-1pt,text=muted] at (0,4) {$4$};
\node[left=-1pt,text=muted] at (0,5) {$5$};
\node[left=-1pt,text=muted] at (0,6) {$6$};
\begin{scope}\clip (-3,-2) rectangle (7,6);
\draw[accent,line width=1.1pt,line join=round] (-3,-2)--(-2.968,-1.968)--(-2.936,-1.936)--(-2.904,-1.904)--(-2.872,-1.872)--(-2.84,-1.84)--(-2.808,-1.808)--(-2.776,-1.776)--(-2.744,-1.744)--(-2.712,-1.712)--(-2.68,-1.68)--(-2.648,-1.648)--(-2.616,-1.616)--(-2.584,-1.584)--(-2.552,-1.552)--(-2.52,-1.52)--(-2.488,-1.488)--(-2.456,-1.456)--(-2.424,-1.424)--(-2.392,-1.392)--(-2.36,-1.36)--(-2.328,-1.328)--(-2.296,-1.296)--(-2.264,-1.264)--(-2.232,-1.232)--(-2.2,-1.2)--(-2.168,-1.168)--(-2.136,-1.136)--(-2.104,-1.104)--(-2.072,-1.072)--(-2.04,-1.04)--(-2.008,-1.008)--(-1.976,-0.976)--(-1.944,-0.944)--(-1.912,-0.912)--(-1.88,-0.88)--(-1.848,-0.848)--(-1.816,-0.816)--(-1.784,-0.784)--(-1.752,-0.752)--(-1.72,-0.72)--(-1.688,-0.688)--(-1.656,-0.656)--(-1.624,-0.624)--(-1.592,-0.592)--(-1.56,-0.56)--(-1.528,-0.528)--(-1.496,-0.496)--(-1.464,-0.464)--(-1.432,-0.432)--(-1.4,-0.4)--(-1.368,-0.368)--(-1.336,-0.336)--(-1.304,-0.304)--(-1.272,-0.272)--(-1.24,-0.24)--(-1.208,-0.208)--(-1.176,-0.176)--(-1.144,-0.144)--(-1.112,-0.112)--(-1.08,-0.08)--(-1.048,-0.048)--(-1.016,-0.016)--(-0.984,0.016)--(-0.952,0.048)--(-0.92,0.08)--(-0.888,0.112)--(-0.856,0.144)--(-0.824,0.176)--(-0.792,0.208)--(-0.76,0.24)--(-0.728,0.272)--(-0.696,0.304)--(-0.664,0.336)--(-0.632,0.368)--(-0.6,0.4)--(-0.568,0.432)--(-0.536,0.464)--(-0.504,0.496)--(-0.472,0.528)--(-0.44,0.56)--(-0.408,0.592)--(-0.376,0.624)--(-0.344,0.656)--(-0.312,0.688)--(-0.28,0.72)--(-0.248,0.752)--(-0.216,0.784)--(-0.184,0.816)--(-0.152,0.848)--(-0.12,0.88)--(-0.088,0.912)--(-0.056,0.944)--(-0.024,0.976)--(0.008,1.008)--(0.04,1.04)--(0.072,1.072)--(0.104,1.104)--(0.136,1.136)--(0.168,1.168)--(0.2,1.2)--(0.232,1.232)--(0.264,1.264)--(0.296,1.296)--(0.328,1.328)--(0.36,1.36)--(0.392,1.392)--(0.424,1.424)--(0.456,1.456)--(0.488,1.488)--(0.52,1.52)--(0.552,1.552)--(0.584,1.584)--(0.616,1.616)--(0.648,1.648)--(0.68,1.68)--(0.712,1.712)--(0.744,1.744)--(0.776,1.776)--(0.808,1.808)--(0.84,1.84)--(0.872,1.872)--(0.904,1.904)--(0.936,1.936)--(0.968,1.968)--(1,2);
\draw[accent,line width=1.1pt,line join=round] (1,0)--(1.048,0.048)--(1.096,0.096)--(1.144,0.144)--(1.192,0.192)--(1.24,0.24)--(1.288,0.288)--(1.336,0.336)--(1.384,0.384)--(1.432,0.432)--(1.48,0.48)--(1.528,0.528)--(1.576,0.576)--(1.624,0.624)--(1.672,0.672)--(1.72,0.72)--(1.768,0.768)--(1.816,0.816)--(1.864,0.864)--(1.912,0.912)--(1.96,0.96)--(2.008,1.008)--(2.056,1.056)--(2.104,1.104)--(2.152,1.152)--(2.2,1.2)--(2.248,1.248)--(2.296,1.296)--(2.344,1.344)--(2.392,1.392)--(2.44,1.44)--(2.488,1.488)--(2.536,1.536)--(2.584,1.584)--(2.632,1.632)--(2.68,1.68)--(2.728,1.728)--(2.776,1.776)--(2.824,1.824)--(2.872,1.872)--(2.92,1.92)--(2.968,1.968)--(3.016,2.016)--(3.064,2.064)--(3.112,2.112)--(3.16,2.16)--(3.208,2.208)--(3.256,2.256)--(3.304,2.304)--(3.352,2.352)--(3.4,2.4)--(3.448,2.448)--(3.496,2.496)--(3.544,2.544)--(3.592,2.592)--(3.64,2.64)--(3.688,2.688)--(3.736,2.736)--(3.784,2.784)--(3.832,2.832)--(3.88,2.88)--(3.928,2.928)--(3.976,2.976)--(4.024,3.024)--(4.072,3.072)--(4.12,3.12)--(4.168,3.168)--(4.216,3.216)--(4.264,3.264)--(4.312,3.312)--(4.36,3.36)--(4.408,3.408)--(4.456,3.456)--(4.504,3.504)--(4.552,3.552)--(4.6,3.6)--(4.648,3.648)--(4.696,3.696)--(4.744,3.744)--(4.792,3.792)--(4.84,3.84)--(4.888,3.888)--(4.936,3.936)--(4.984,3.984)--(5.032,4.032)--(5.08,4.08)--(5.128,4.128)--(5.176,4.176)--(5.224,4.224)--(5.272,4.272)--(5.32,4.32)--(5.368,4.368)--(5.416,4.416)--(5.464,4.464)--(5.512,4.512)--(5.56,4.56)--(5.608,4.608)--(5.656,4.656)--(5.704,4.704)--(5.752,4.752)--(5.8,4.8)--(5.848,4.848)--(5.896,4.896)--(5.944,4.944)--(5.992,4.992)--(6.04,5.04)--(6.088,5.088)--(6.136,5.136)--(6.184,5.184)--(6.232,5.232)--(6.28,5.28)--(6.328,5.328)--(6.376,5.376)--(6.424,5.424)--(6.472,5.472)--(6.52,5.52)--(6.568,5.568)--(6.616,5.616)--(6.664,5.664)--(6.712,5.712)--(6.76,5.76)--(6.808,5.808)--(6.856,5.856)--(6.904,5.904)--(6.952,5.952)--(7,6);
\filldraw[fill=accent,draw=accent,line width=.8pt] (1,2) circle (1.5pt);
\filldraw[fill=white,draw=accent,line width=.8pt] (1,0) circle (1.5pt);
\filldraw[fill=white,draw=accent,line width=.8pt] (4,3) circle (1.5pt);
\filldraw[fill=accent,draw=accent,line width=.8pt] (4,5) circle (1.5pt);
\end{scope}\end{tikzpicture}\end{center}


\ex{2} \begin{enumerate}[label=\alph*),leftmargin=*,itemsep=2pt,topsep=2pt]\item 40 min: \rest{1,50 €}; \ 1 h 20 min (dos horas empezadas): \rest{2,50 €}; \ 3 h justas: \rest{3,50 €}.\item Es una función «a escalones»:

\begin{center}\begin{tikzpicture}[x=0.76cm,y=0.694cm,baseline={(current bounding box.center)},font=\scriptsize]
\draw[muted!25,line width=.3pt] (0,-0.8)--(0,5.5);
\draw[muted!25,line width=.3pt] (1,-0.8)--(1,5.5);
\draw[muted!25,line width=.3pt] (2,-0.8)--(2,5.5);
\draw[muted!25,line width=.3pt] (3,-0.8)--(3,5.5);
\draw[muted!25,line width=.3pt] (4,-0.8)--(4,5.5);
\draw[muted!25,line width=.3pt] (-0.4,0)--(4.6,0);
\draw[muted!25,line width=.3pt] (-0.4,1)--(4.6,1);
\draw[muted!25,line width=.3pt] (-0.4,2)--(4.6,2);
\draw[muted!25,line width=.3pt] (-0.4,3)--(4.6,3);
\draw[muted!25,line width=.3pt] (-0.4,4)--(4.6,4);
\draw[muted!25,line width=.3pt] (-0.4,5)--(4.6,5);
\draw[-latex,muted] (-0.4,0)--(4.6,0);
\draw[-latex,muted] (0,-0.8)--(0,5.5);
\node[below=-1pt,text=muted] at (1,0) {$1$};
\node[below=-1pt,text=muted] at (2,0) {$2$};
\node[below=-1pt,text=muted] at (3,0) {$3$};
\node[below=-1pt,text=muted] at (4,0) {$4$};
\node[left=-1pt,text=muted] at (0,1) {$1$};
\node[left=-1pt,text=muted] at (0,2) {$2$};
\node[left=-1pt,text=muted] at (0,3) {$3$};
\node[left=-1pt,text=muted] at (0,4) {$4$};
\node[left=-1pt,text=muted] at (0,5) {$5$};
\begin{scope}\clip (-0.4,-0.8) rectangle (4.6,5.5);
\draw[accent,line width=1.1pt,line join=round] (0,1.5)--(1,1.5);
\draw[accent,line width=1.1pt,line join=round] (1,2.5)--(2,2.5);
\draw[accent,line width=1.1pt,line join=round] (2,3.5)--(3,3.5);
\draw[accent,line width=1.1pt,line join=round] (3,4.5)--(4,4.5);
\filldraw[fill=white,draw=accent,line width=.8pt] (0,1.5) circle (1.5pt);
\filldraw[fill=accent,draw=accent,line width=.8pt] (1,1.5) circle (1.5pt);
\filldraw[fill=white,draw=accent,line width=.8pt] (1,2.5) circle (1.5pt);
\filldraw[fill=accent,draw=accent,line width=.8pt] (2,2.5) circle (1.5pt);
\filldraw[fill=white,draw=accent,line width=.8pt] (2,3.5) circle (1.5pt);
\filldraw[fill=accent,draw=accent,line width=.8pt] (3,3.5) circle (1.5pt);
\filldraw[fill=white,draw=accent,line width=.8pt] (3,4.5) circle (1.5pt);
\filldraw[fill=accent,draw=accent,line width=.8pt] (4,4.5) circle (1.5pt);
\end{scope}\end{tikzpicture}\end{center}
\item No: tiene \rest{discontinuidades de salto finito} en $x=1$, $x=2$ y $x=3$ (el precio sube 1 € de golpe). $R=\res{\{1{,}5;\ 2{,}5;\ 3{,}5;\ 4{,}5\}}$.\end{enumerate}

\ex{3} $D=\res{[-6,\,7]}$. \ $x=-3$: $\res{\text{evitable}}$ (los dos lados llegan a 2, pero $f(-3)=4$). \ $x=1$: $\res{\text{salto infinito}}$ (por la izquierda llega a 0 y por la derecha se va a $+\infty$). \ $x=3$: $\res{\text{salto finito}}$ (por la izquierda llega a 1 y $f(3)=2$).


% =====================================================
\seccio{5}{Crecimiento y decrecimiento. Máximos y mínimos}

\exf{1} $D=\res{[-5,\,6]}$, $R=\res{[-3,\,4]}$. Creciente en $\res{(-5,\,-3)\cup(-1,\,1)\cup(3,\,6)}$; decreciente en $\res{(-3,\,-1)\cup(1,\,3)}$. Máximos: $\res{(-3,\,3)}$ relativo y $\res{(1,\,4)}$ absoluto. Mínimos: $\res{(-1,\,1)}$ relativo y $\res{(3,\,-3)}$ absoluto.

\ex{1} Creciente en $\res{(-\infty,\,2)}$ y decreciente en $\res{(2,\,+\infty)}$; máximo absoluto en $\res{(2,\,4)}$: $f(2)=-4+8=4$. No tiene mínimos. Cortes: $-x^2+4x=0\Rightarrow x(4-x)=0\Rightarrow(0,\,0)$ y $(4,\,0)$.

\ex{1} Creciente en $\res{(-\infty,\,-1)\cup(1,\,+\infty)}$; decreciente en $\res{(-1,\,1)}$. Máximo relativo $\res{(-1,\,2)}$ y mínimo relativo $\res{(1,\,-2)}$. No tiene extremos absolutos: por la derecha sube sin parar y por la izquierda baja sin parar.

\ex{2} Por ejemplo (hay muchas), la parábola $f(x)=\frac12(x-1)^2-2$:

\begin{center}\begin{tikzpicture}[x=0.463cm,y=0.463cm,baseline={(current bounding box.center)},font=\scriptsize]
\draw[muted!25,line width=.3pt] (-3,-3)--(-3,4);
\draw[muted!25,line width=.3pt] (-2,-3)--(-2,4);
\draw[muted!25,line width=.3pt] (-1,-3)--(-1,4);
\draw[muted!25,line width=.3pt] (0,-3)--(0,4);
\draw[muted!25,line width=.3pt] (1,-3)--(1,4);
\draw[muted!25,line width=.3pt] (2,-3)--(2,4);
\draw[muted!25,line width=.3pt] (3,-3)--(3,4);
\draw[muted!25,line width=.3pt] (4,-3)--(4,4);
\draw[muted!25,line width=.3pt] (5,-3)--(5,4);
\draw[muted!25,line width=.3pt] (-3,-3)--(5,-3);
\draw[muted!25,line width=.3pt] (-3,-2)--(5,-2);
\draw[muted!25,line width=.3pt] (-3,-1)--(5,-1);
\draw[muted!25,line width=.3pt] (-3,0)--(5,0);
\draw[muted!25,line width=.3pt] (-3,1)--(5,1);
\draw[muted!25,line width=.3pt] (-3,2)--(5,2);
\draw[muted!25,line width=.3pt] (-3,3)--(5,3);
\draw[muted!25,line width=.3pt] (-3,4)--(5,4);
\draw[-latex,muted] (-3,0)--(5,0);
\draw[-latex,muted] (0,-3)--(0,4);
\node[below=-1pt,text=muted] at (-3,0) {$-3$};
\node[below=-1pt,text=muted] at (-2,0) {$-2$};
\node[below=-1pt,text=muted] at (-1,0) {$-1$};
\node[below=-1pt,text=muted] at (1,0) {$1$};
\node[below=-1pt,text=muted] at (2,0) {$2$};
\node[below=-1pt,text=muted] at (3,0) {$3$};
\node[below=-1pt,text=muted] at (4,0) {$4$};
\node[below=-1pt,text=muted] at (5,0) {$5$};
\node[left=-1pt,text=muted] at (0,-3) {$-3$};
\node[left=-1pt,text=muted] at (0,-2) {$-2$};
\node[left=-1pt,text=muted] at (0,-1) {$-1$};
\node[left=-1pt,text=muted] at (0,1) {$1$};
\node[left=-1pt,text=muted] at (0,2) {$2$};
\node[left=-1pt,text=muted] at (0,3) {$3$};
\node[left=-1pt,text=muted] at (0,4) {$4$};
\begin{scope}\clip (-3,-3) rectangle (5,4);
\draw[accent,line width=1.1pt,line join=round] (-3,6)--(-2.936,5.746)--(-2.872,5.496)--(-2.808,5.25)--(-2.744,5.009)--(-2.68,4.771)--(-2.616,4.538)--(-2.552,4.308)--(-2.488,4.083)--(-2.424,3.862)--(-2.36,3.645)--(-2.296,3.432)--(-2.232,3.223)--(-2.168,3.018)--(-2.104,2.817)--(-2.04,2.621)--(-1.976,2.428)--(-1.912,2.24)--(-1.848,2.056)--(-1.784,1.875)--(-1.72,1.699)--(-1.656,1.527)--(-1.592,1.359)--(-1.528,1.195)--(-1.464,1.036)--(-1.4,0.88)--(-1.336,0.728)--(-1.272,0.581)--(-1.208,0.438)--(-1.144,0.298)--(-1.08,0.163)--(-1.016,0.032)--(-0.952,-0.095)--(-0.888,-0.218)--(-0.824,-0.337)--(-0.76,-0.451)--(-0.696,-0.562)--(-0.632,-0.668)--(-0.568,-0.771)--(-0.504,-0.869)--(-0.44,-0.963)--(-0.376,-1.053)--(-0.312,-1.139)--(-0.248,-1.221)--(-0.184,-1.299)--(-0.12,-1.373)--(-0.056,-1.442)--(0.008,-1.508)--(0.072,-1.569)--(0.136,-1.627)--(0.2,-1.68)--(0.264,-1.729)--(0.328,-1.774)--(0.392,-1.815)--(0.456,-1.852)--(0.52,-1.885)--(0.584,-1.913)--(0.648,-1.938)--(0.712,-1.959)--(0.776,-1.975)--(0.84,-1.987)--(0.904,-1.995)--(0.968,-1.999)--(1.032,-1.999)--(1.096,-1.995)--(1.16,-1.987)--(1.224,-1.975)--(1.288,-1.959)--(1.352,-1.938)--(1.416,-1.913)--(1.48,-1.885)--(1.544,-1.852)--(1.608,-1.815)--(1.672,-1.774)--(1.736,-1.729)--(1.8,-1.68)--(1.864,-1.627)--(1.928,-1.569)--(1.992,-1.508)--(2.056,-1.442)--(2.12,-1.373)--(2.184,-1.299)--(2.248,-1.221)--(2.312,-1.139)--(2.376,-1.053)--(2.44,-0.963)--(2.504,-0.869)--(2.568,-0.771)--(2.632,-0.668)--(2.696,-0.562)--(2.76,-0.451)--(2.824,-0.337)--(2.888,-0.218)--(2.952,-0.095)--(3.016,0.032)--(3.08,0.163)--(3.144,0.298)--(3.208,0.438)--(3.272,0.581)--(3.336,0.728)--(3.4,0.88)--(3.464,1.036)--(3.528,1.195)--(3.592,1.359)--(3.656,1.527)--(3.72,1.699)--(3.784,1.875)--(3.848,2.056)--(3.912,2.24)--(3.976,2.428)--(4.04,2.621)--(4.104,2.817)--(4.168,3.018)--(4.232,3.223)--(4.296,3.432)--(4.36,3.645)--(4.424,3.862)--(4.488,4.083)--(4.552,4.308)--(4.616,4.538)--(4.68,4.771)--(4.744,5.009)--(4.808,5.25)--(4.872,5.496)--(4.936,5.746)--(5,6);
\filldraw[fill=blue!70!black,draw=blue!70!black,line width=.8pt] (1,-2) circle (1.5pt);
\filldraw[fill=blue!70!black,draw=blue!70!black,line width=.8pt] (-1,0) circle (1.5pt);
\filldraw[fill=blue!70!black,draw=blue!70!black,line width=.8pt] (3,0) circle (1.5pt);
\end{scope}\end{tikzpicture}\end{center}


\ex{2} 

\begin{center}\begin{tikzpicture}[x=0.231cm,y=0.231cm,baseline={(current bounding box.center)},font=\scriptsize]
\draw[muted!25,line width=.3pt] (0,-1)--(0,23);
\draw[muted!25,line width=.3pt] (2,-1)--(2,23);
\draw[muted!25,line width=.3pt] (4,-1)--(4,23);
\draw[muted!25,line width=.3pt] (6,-1)--(6,23);
\draw[muted!25,line width=.3pt] (8,-1)--(8,23);
\draw[muted!25,line width=.3pt] (10,-1)--(10,23);
\draw[muted!25,line width=.3pt] (12,-1)--(12,23);
\draw[muted!25,line width=.3pt] (14,-1)--(14,23);
\draw[muted!25,line width=.3pt] (16,-1)--(16,23);
\draw[muted!25,line width=.3pt] (18,-1)--(18,23);
\draw[muted!25,line width=.3pt] (20,-1)--(20,23);
\draw[muted!25,line width=.3pt] (22,-1)--(22,23);
\draw[muted!25,line width=.3pt] (24,-1)--(24,23);
\draw[muted!25,line width=.3pt] (-1,0)--(25,0);
\draw[muted!25,line width=.3pt] (-1,2)--(25,2);
\draw[muted!25,line width=.3pt] (-1,4)--(25,4);
\draw[muted!25,line width=.3pt] (-1,6)--(25,6);
\draw[muted!25,line width=.3pt] (-1,8)--(25,8);
\draw[muted!25,line width=.3pt] (-1,10)--(25,10);
\draw[muted!25,line width=.3pt] (-1,12)--(25,12);
\draw[muted!25,line width=.3pt] (-1,14)--(25,14);
\draw[muted!25,line width=.3pt] (-1,16)--(25,16);
\draw[muted!25,line width=.3pt] (-1,18)--(25,18);
\draw[muted!25,line width=.3pt] (-1,20)--(25,20);
\draw[muted!25,line width=.3pt] (-1,22)--(25,22);
\draw[-latex,muted] (-1,0)--(25,0);
\draw[-latex,muted] (0,-1)--(0,23);
\node[below=-1pt,text=muted] at (4,0) {$4$};
\node[below=-1pt,text=muted] at (8,0) {$8$};
\node[below=-1pt,text=muted] at (12,0) {$12$};
\node[below=-1pt,text=muted] at (16,0) {$16$};
\node[below=-1pt,text=muted] at (20,0) {$20$};
\node[below=-1pt,text=muted] at (24,0) {$24$};
\node[left=-1pt,text=muted] at (0,4) {$4$};
\node[left=-1pt,text=muted] at (0,8) {$8$};
\node[left=-1pt,text=muted] at (0,12) {$12$};
\node[left=-1pt,text=muted] at (0,16) {$16$};
\node[left=-1pt,text=muted] at (0,20) {$20$};
\begin{scope}\clip (-1,-1) rectangle (25,23);
\draw[accent,line width=1.1pt,line join=round] (0,8)--(2,6);
\draw[accent,line width=1.1pt,line join=round] (2,6)--(4,5);
\draw[accent,line width=1.1pt,line join=round] (4,5)--(6,7);
\draw[accent,line width=1.1pt,line join=round] (6,7)--(8,12);
\draw[accent,line width=1.1pt,line join=round] (8,12)--(10,16);
\draw[accent,line width=1.1pt,line join=round] (10,16)--(12,19);
\draw[accent,line width=1.1pt,line join=round] (12,19)--(14,21);
\draw[accent,line width=1.1pt,line join=round] (14,21)--(16,20);
\draw[accent,line width=1.1pt,line join=round] (16,20)--(18,16);
\draw[accent,line width=1.1pt,line join=round] (18,16)--(20,13);
\draw[accent,line width=1.1pt,line join=round] (20,13)--(22,11);
\draw[accent,line width=1.1pt,line join=round] (22,11)--(24,9);
\filldraw[fill=accent,draw=accent,line width=.8pt] (0,8) circle (1.5pt);
\filldraw[fill=accent,draw=accent,line width=.8pt] (2,6) circle (1.5pt);
\filldraw[fill=accent,draw=accent,line width=.8pt] (4,5) circle (1.5pt);
\filldraw[fill=accent,draw=accent,line width=.8pt] (6,7) circle (1.5pt);
\filldraw[fill=accent,draw=accent,line width=.8pt] (8,12) circle (1.5pt);
\filldraw[fill=accent,draw=accent,line width=.8pt] (10,16) circle (1.5pt);
\filldraw[fill=accent,draw=accent,line width=.8pt] (12,19) circle (1.5pt);
\filldraw[fill=accent,draw=accent,line width=.8pt] (14,21) circle (1.5pt);
\filldraw[fill=accent,draw=accent,line width=.8pt] (16,20) circle (1.5pt);
\filldraw[fill=accent,draw=accent,line width=.8pt] (18,16) circle (1.5pt);
\filldraw[fill=accent,draw=accent,line width=.8pt] (20,13) circle (1.5pt);
\filldraw[fill=accent,draw=accent,line width=.8pt] (22,11) circle (1.5pt);
\filldraw[fill=accent,draw=accent,line width=.8pt] (24,9) circle (1.5pt);
\end{scope}\end{tikzpicture}\end{center}
Baja de las 0 h a las \rest{4 h}, sube de las 4 h a las \rest{14 h} y vuelve a bajar hasta las 24 h. Mínima: \rest{5 °C a las 4 h}; máxima: \rest{21 °C a las 14 h}. Unimos los puntos porque la temperatura cambia continuamente (los datos son aproximados entre las medidas).

\ex{2} \begin{enumerate}[label=\alph*),leftmargin=*,itemsep=2pt,topsep=2pt]\item $D=\res{[0,\,12]}$: el vuelo dura 12 min. $R=\res{[0,\,50]}$: la altura va de 0 m a 50 m.\item Sube en $\res{(0,\,2)\cup(5,\,7)}$, se mantiene en $\res{(2,\,5)}$ y baja en $\res{(7,\,12)}$.\item Máximo absoluto $\res{(7,\,50)}$: 50 m a los 7 min.\item Desde el minuto 5 hasta el minuto 9 (de 7 a 10 baja 10 m cada minuto): \rest{4 minutos}.\end{enumerate}


% =====================================================
\seccio{6}{Periodicidad}

\exf{1} \begin{tasks}(1)
  \task $\res{T=3}$
  \task $f(10)=f(1+3\cdot3)=f(1)=\res{2}$; \ $f(-1)=f(-1+3)=f(2)=\res{3}$; \ $f(30)=f(0+10\cdot3)=f(0)=\res{1}$; \ $f(47)=f(2+15\cdot3)=f(2)=\res{3}$
\end{tasks}

\ex{2} $f(11)=f(1+2\cdot5)=\res{4}$; \ $f(-2)=f(-2+5)=f(3)=\res{-2}$; \ $f(23)=f(3+4\cdot5)=\res{-2}$; \ $f(101)=f(1+20\cdot5)=\res{4}$

\ex{2} \begin{enumerate}[label=\alph*),leftmargin=*,itemsep=2pt,topsep=2pt]\item $\res{T=60\text{ s}}$. Altura mínima 2 m y máxima $2+20=22$ m: $R=\res{[2,\,22]}$.\item $150=2\cdot60+30$: como a los 30 s (media vuelta), está arriba del todo: \rest{22 m}. \ $240=4\cdot60$: como al inicio: \rest{2 m}.\item Da $600:60=10$ vueltas: \rest{10 veces}.\end{enumerate}

\ex{3} 

\begin{center}\begin{tikzpicture}[x=0.571cm,y=1.24cm,baseline={(current bounding box.center)},font=\scriptsize]
\draw[muted!25,line width=.3pt] (-5,-1)--(-5,3);
\draw[muted!25,line width=.3pt] (-4,-1)--(-4,3);
\draw[muted!25,line width=.3pt] (-3,-1)--(-3,3);
\draw[muted!25,line width=.3pt] (-2,-1)--(-2,3);
\draw[muted!25,line width=.3pt] (-1,-1)--(-1,3);
\draw[muted!25,line width=.3pt] (0,-1)--(0,3);
\draw[muted!25,line width=.3pt] (1,-1)--(1,3);
\draw[muted!25,line width=.3pt] (2,-1)--(2,3);
\draw[muted!25,line width=.3pt] (3,-1)--(3,3);
\draw[muted!25,line width=.3pt] (4,-1)--(4,3);
\draw[muted!25,line width=.3pt] (5,-1)--(5,3);
\draw[muted!25,line width=.3pt] (6,-1)--(6,3);
\draw[muted!25,line width=.3pt] (7,-1)--(7,3);
\draw[muted!25,line width=.3pt] (8,-1)--(8,3);
\draw[muted!25,line width=.3pt] (9,-1)--(9,3);
\draw[muted!25,line width=.3pt] (-5,-1)--(9,-1);
\draw[muted!25,line width=.3pt] (-5,0)--(9,0);
\draw[muted!25,line width=.3pt] (-5,1)--(9,1);
\draw[muted!25,line width=.3pt] (-5,2)--(9,2);
\draw[muted!25,line width=.3pt] (-5,3)--(9,3);
\draw[-latex,muted] (-5,0)--(9,0);
\draw[-latex,muted] (0,-1)--(0,3);
\node[below=-1pt,text=muted] at (-5,0) {$-5$};
\node[below=-1pt,text=muted] at (-4,0) {$-4$};
\node[below=-1pt,text=muted] at (-3,0) {$-3$};
\node[below=-1pt,text=muted] at (-2,0) {$-2$};
\node[below=-1pt,text=muted] at (-1,0) {$-1$};
\node[below=-1pt,text=muted] at (1,0) {$1$};
\node[below=-1pt,text=muted] at (2,0) {$2$};
\node[below=-1pt,text=muted] at (3,0) {$3$};
\node[below=-1pt,text=muted] at (4,0) {$4$};
\node[below=-1pt,text=muted] at (5,0) {$5$};
\node[below=-1pt,text=muted] at (6,0) {$6$};
\node[below=-1pt,text=muted] at (7,0) {$7$};
\node[below=-1pt,text=muted] at (8,0) {$8$};
\node[below=-1pt,text=muted] at (9,0) {$9$};
\node[left=-1pt,text=muted] at (0,-1) {$-1$};
\node[left=-1pt,text=muted] at (0,1) {$1$};
\node[left=-1pt,text=muted] at (0,2) {$2$};
\node[left=-1pt,text=muted] at (0,3) {$3$};
\begin{scope}\clip (-5,-1) rectangle (9,3);
\draw[accent,line width=1.1pt,line join=round] (-4,2)--(-3.9,1.9)--(-3.8,1.8)--(-3.7,1.7)--(-3.6,1.6)--(-3.5,1.5)--(-3.4,1.4)--(-3.3,1.3)--(-3.2,1.2)--(-3.1,1.1)--(-3,1)--(-2.9,0.9)--(-2.8,0.8)--(-2.7,0.7)--(-2.6,0.6)--(-2.5,0.5)--(-2.4,0.4)--(-2.3,0.3)--(-2.2,0.2)--(-2.1,0.1)--(-2,0)--(-1.9,0.1)--(-1.8,0.2)--(-1.7,0.3)--(-1.6,0.4)--(-1.5,0.5)--(-1.4,0.6)--(-1.3,0.7)--(-1.2,0.8)--(-1.1,0.9)--(-1,1)--(-0.9,1.1)--(-0.8,1.2)--(-0.7,1.3)--(-0.6,1.4)--(-0.5,1.5)--(-0.4,1.6)--(-0.3,1.7)--(-0.2,1.8)--(-0.1,1.9)--(0,2)--(0.1,1.9)--(0.2,1.8)--(0.3,1.7)--(0.4,1.6)--(0.5,1.5)--(0.6,1.4)--(0.7,1.3)--(0.8,1.2)--(0.9,1.1)--(1,1)--(1.1,0.9)--(1.2,0.8)--(1.3,0.7)--(1.4,0.6)--(1.5,0.5)--(1.6,0.4)--(1.7,0.3)--(1.8,0.2)--(1.9,0.1)--(2,0)--(2.1,0.1)--(2.2,0.2)--(2.3,0.3)--(2.4,0.4)--(2.5,0.5)--(2.6,0.6)--(2.7,0.7)--(2.8,0.8)--(2.9,0.9)--(3,1)--(3.1,1.1)--(3.2,1.2)--(3.3,1.3)--(3.4,1.4)--(3.5,1.5)--(3.6,1.6)--(3.7,1.7)--(3.8,1.8)--(3.9,1.9)--(4,2)--(4.1,1.9)--(4.2,1.8)--(4.3,1.7)--(4.4,1.6)--(4.5,1.5)--(4.6,1.4)--(4.7,1.3)--(4.8,1.2)--(4.9,1.1)--(5,1)--(5.1,0.9)--(5.2,0.8)--(5.3,0.7)--(5.4,0.6)--(5.5,0.5)--(5.6,0.4)--(5.7,0.3)--(5.8,0.2)--(5.9,0.1)--(6,0)--(6.1,0.1)--(6.2,0.2)--(6.3,0.3)--(6.4,0.4)--(6.5,0.5)--(6.6,0.6)--(6.7,0.7)--(6.8,0.8)--(6.9,0.9)--(7,1)--(7.1,1.1)--(7.2,1.2)--(7.3,1.3)--(7.4,1.4)--(7.5,1.5)--(7.6,1.6)--(7.7,1.7)--(7.8,1.8)--(7.9,1.9)--(8,2);
\draw[blue!70!black,line width=1.1pt,line join=round] (0,2)--(0.03,1.97)--(0.06,1.94)--(0.09,1.91)--(0.12,1.88)--(0.15,1.85)--(0.18,1.82)--(0.21,1.79)--(0.24,1.76)--(0.27,1.73)--(0.3,1.7)--(0.33,1.67)--(0.36,1.64)--(0.39,1.61)--(0.42,1.58)--(0.45,1.55)--(0.48,1.52)--(0.51,1.49)--(0.54,1.46)--(0.57,1.43)--(0.6,1.4)--(0.63,1.37)--(0.66,1.34)--(0.69,1.31)--(0.72,1.28)--(0.75,1.25)--(0.78,1.22)--(0.81,1.19)--(0.84,1.16)--(0.87,1.13)--(0.9,1.1)--(0.93,1.07)--(0.96,1.04)--(0.99,1.01)--(1.02,0.98)--(1.05,0.95)--(1.08,0.92)--(1.11,0.89)--(1.14,0.86)--(1.17,0.83)--(1.2,0.8)--(1.23,0.77)--(1.26,0.74)--(1.29,0.71)--(1.32,0.68)--(1.35,0.65)--(1.38,0.62)--(1.41,0.59)--(1.44,0.56)--(1.47,0.53)--(1.5,0.5)--(1.53,0.47)--(1.56,0.44)--(1.59,0.41)--(1.62,0.38)--(1.65,0.35)--(1.68,0.32)--(1.71,0.29)--(1.74,0.26)--(1.77,0.23)--(1.8,0.2)--(1.83,0.17)--(1.86,0.14)--(1.89,0.11)--(1.92,0.08)--(1.95,0.05)--(1.98,0.02)--(2.01,0.01)--(2.04,0.04)--(2.07,0.07)--(2.1,0.1)--(2.13,0.13)--(2.16,0.16)--(2.19,0.19)--(2.22,0.22)--(2.25,0.25)--(2.28,0.28)--(2.31,0.31)--(2.34,0.34)--(2.37,0.37)--(2.4,0.4)--(2.43,0.43)--(2.46,0.46)--(2.49,0.49)--(2.52,0.52)--(2.55,0.55)--(2.58,0.58)--(2.61,0.61)--(2.64,0.64)--(2.67,0.67)--(2.7,0.7)--(2.73,0.73)--(2.76,0.76)--(2.79,0.79)--(2.82,0.82)--(2.85,0.85)--(2.88,0.88)--(2.91,0.91)--(2.94,0.94)--(2.97,0.97)--(3,1)--(3.03,1.03)--(3.06,1.06)--(3.09,1.09)--(3.12,1.12)--(3.15,1.15)--(3.18,1.18)--(3.21,1.21)--(3.24,1.24)--(3.27,1.27)--(3.3,1.3)--(3.33,1.33)--(3.36,1.36)--(3.39,1.39)--(3.42,1.42)--(3.45,1.45)--(3.48,1.48)--(3.51,1.51)--(3.54,1.54)--(3.57,1.57)--(3.6,1.6)--(3.63,1.63)--(3.66,1.66)--(3.69,1.69)--(3.72,1.72)--(3.75,1.75)--(3.78,1.78)--(3.81,1.81)--(3.84,1.84)--(3.87,1.87)--(3.9,1.9)--(3.93,1.93)--(3.96,1.96)--(3.99,1.99)--(4,2);
\end{scope}\end{tikzpicture}\end{center}
$f(13)=f(1+3\cdot4)=f(1)=|1-2|=\res{1}$

\ex{1} \begin{tasks}(1)
  \task \rest{Sí}: se repite cada hora, $T=\res{60\text{ min}}$.
  \task \rest{No}: crece hasta la edad adulta y no se repite.
  \task \rest{Sí}: $T=40+5+35=\res{80\text{ s}}$.
  \task \rest{No}: sube y baja, pero no se repite de manera regular.
\end{tasks}


% =====================================================
\seccio{7}{Interpretación de gráficos}

\exf{1} \begin{enumerate}[label=\alph*),leftmargin=*,itemsep=2pt,topsep=2pt]\item Independiente: tiempo (min); dependiente: volumen (L). $D=\res{[0,\,25]}$: todo dura 25 minutos. $R=\res{[0,\,200]}$: el depósito tiene entre 0 y 200 L.\item En $(0,\,10)$ \rest{se llena a 20 L/min}; en $(10,\,15)$ \rest{se mantiene lleno} (200 L); en $(15,\,25)$ \rest{se vacía a 20 L/min}.\item A los 12 min: \rest{200 L}. Hay 100 L \rest{a los 5 min} y \rest{a los 20 min}.\end{enumerate}

\ex{2} \begin{enumerate}[label=\alph*),leftmargin=*,itemsep=2pt,topsep=2pt]\item $D=\res{[0,\,5]}$ y $R=\res{[0,\,40]}$: la salida dura 5 h y llega a 40 km de casa (máximo absoluto $(3,\,40)$).\item Creciente en $\res{(0,\,1)\cup(2,\,3)}$, constante en $\res{(1,\,2)}$ (está parada) y decreciente en $\res{(3,\,5)}$ (vuelve).\item Hace 40 km en 2 h: \rest{20 km/h}.\item A la \rest{media hora} (ida) y a las \rest{4 horas y media} (vuelta).\end{enumerate}

\ex{1} \begin{tasks}(4)
  \task gráfica \rest{2}
  \task gráfica \rest{4}
  \task gráfica \rest{3}
  \task gráfica \rest{1}
\end{tasks}

\ex{2} $D=\res{[-5,\,6]}$, $R=\res{[-2,\,4]}$. Cortes: $OX$: $\res{(-5,\,0),\ (-1,\,0),\ (1,\,0),\ (5,\,0)}$; $OY$: $\res{(0,\,-1)}$. Discontinua en $x=0$ (\rest{salto finito}: por la izquierda llega a $-1=f(0)$ y por la derecha empieza en $-2$). Creciente en $\res{(-5,\,-3)\cup(0,\,3)}$; decreciente en $\res{(-3,\,0)\cup(3,\,6)}$. Máximo relativo $\res{(-3,\,2)}$ y máximo absoluto $\res{(3,\,4)}$; mínimo absoluto $\res{(6,\,-2)}$.

\ex{3} En 15 min a 6 km/h hace $6\cdot\frac14=1{,}5$ km. Puntos: $(0,\,0)$, $(15;\ 1{,}5)$, $(25;\ 1{,}5)$, $(40;\ 4{,}5)$, $(70,\,0)$ (el tiempo del polideportivo no cuenta). Velocidades: \rest{6 km/h}, \rest{0} (parado), $3:\frac14=$ \rest{12 km/h} y $4{,}5:\frac12=$ \rest{9 km/h}.

\begin{center}\begin{tikzpicture}[x=0.062cm,y=0.517cm,baseline={(current bounding box.center)},font=\scriptsize]
\draw[muted!25,line width=.3pt] (0,-0.5)--(0,5.5);
\draw[muted!25,line width=.3pt] (5,-0.5)--(5,5.5);
\draw[muted!25,line width=.3pt] (10,-0.5)--(10,5.5);
\draw[muted!25,line width=.3pt] (15,-0.5)--(15,5.5);
\draw[muted!25,line width=.3pt] (20,-0.5)--(20,5.5);
\draw[muted!25,line width=.3pt] (25,-0.5)--(25,5.5);
\draw[muted!25,line width=.3pt] (30,-0.5)--(30,5.5);
\draw[muted!25,line width=.3pt] (35,-0.5)--(35,5.5);
\draw[muted!25,line width=.3pt] (40,-0.5)--(40,5.5);
\draw[muted!25,line width=.3pt] (45,-0.5)--(45,5.5);
\draw[muted!25,line width=.3pt] (50,-0.5)--(50,5.5);
\draw[muted!25,line width=.3pt] (55,-0.5)--(55,5.5);
\draw[muted!25,line width=.3pt] (60,-0.5)--(60,5.5);
\draw[muted!25,line width=.3pt] (65,-0.5)--(65,5.5);
\draw[muted!25,line width=.3pt] (70,-0.5)--(70,5.5);
\draw[muted!25,line width=.3pt] (75,-0.5)--(75,5.5);
\draw[muted!25,line width=.3pt] (-4,-0.5)--(76,-0.5);
\draw[muted!25,line width=.3pt] (-4,0)--(76,0);
\draw[muted!25,line width=.3pt] (-4,0.5)--(76,0.5);
\draw[muted!25,line width=.3pt] (-4,1)--(76,1);
\draw[muted!25,line width=.3pt] (-4,1.5)--(76,1.5);
\draw[muted!25,line width=.3pt] (-4,2)--(76,2);
\draw[muted!25,line width=.3pt] (-4,2.5)--(76,2.5);
\draw[muted!25,line width=.3pt] (-4,3)--(76,3);
\draw[muted!25,line width=.3pt] (-4,3.5)--(76,3.5);
\draw[muted!25,line width=.3pt] (-4,4)--(76,4);
\draw[muted!25,line width=.3pt] (-4,4.5)--(76,4.5);
\draw[muted!25,line width=.3pt] (-4,5)--(76,5);
\draw[muted!25,line width=.3pt] (-4,5.5)--(76,5.5);
\draw[-latex,muted] (-4,0)--(76,0);
\draw[-latex,muted] (0,-0.5)--(0,5.5);
\node[below=-1pt,text=muted] at (10,0) {$10$};
\node[below=-1pt,text=muted] at (20,0) {$20$};
\node[below=-1pt,text=muted] at (30,0) {$30$};
\node[below=-1pt,text=muted] at (40,0) {$40$};
\node[below=-1pt,text=muted] at (50,0) {$50$};
\node[below=-1pt,text=muted] at (60,0) {$60$};
\node[below=-1pt,text=muted] at (70,0) {$70$};
\node[left=-1pt,text=muted] at (0,1) {$1$};
\node[left=-1pt,text=muted] at (0,2) {$2$};
\node[left=-1pt,text=muted] at (0,3) {$3$};
\node[left=-1pt,text=muted] at (0,4) {$4$};
\node[left=-1pt,text=muted] at (0,5) {$5$};
\begin{scope}\clip (-4,-0.5) rectangle (76,5.5);
\draw[accent,line width=1.1pt,line join=round] (0,0)--(15,1.5);
\draw[accent,line width=1.1pt,line join=round] (15,1.5)--(25,1.5);
\draw[accent,line width=1.1pt,line join=round] (25,1.5)--(40,4.5);
\draw[accent,line width=1.1pt,line join=round] (40,4.5)--(70,0);
\filldraw[fill=accent,draw=accent,line width=.8pt] (0,0) circle (1.5pt);
\filldraw[fill=accent,draw=accent,line width=.8pt] (15,1.5) circle (1.5pt);
\filldraw[fill=accent,draw=accent,line width=.8pt] (25,1.5) circle (1.5pt);
\filldraw[fill=accent,draw=accent,line width=.8pt] (40,4.5) circle (1.5pt);
\filldraw[fill=accent,draw=accent,line width=.8pt] (70,0) circle (1.5pt);
\end{scope}\end{tikzpicture}\end{center}


\ex{1} \begin{enumerate}[label=\alph*),leftmargin=*,itemsep=2pt,topsep=2pt]\item Acelera en $\res{(0,\,2)\cup(8,\,10)}$; velocidad constante en $\res{(2,\,6)}$; frena en $\res{(6,\,7)}$.\item Máxima: \rest{60 km/h}, de 2 a 6 min. Está parado \rest{del minuto 7 al 8}: quizá un semáforo en rojo.\item $\res{20\text{ km/h}}$\end{enumerate}

\ex{2} \begin{enumerate}[label=\alph*),leftmargin=*,itemsep=2pt,topsep=2pt]\item $20t-5t^2=0\Rightarrow 5t(4-t)=0\Rightarrow t=0,\ t=4$: vuelve al suelo a los \rest{4 s}; $D=\res{[0,\,4]}$.\item Máximo $\res{(2,\,20)}$: \rest{20 m a los 2 s} (a la mitad del vuelo, por simetría).\item $20t-5t^2=15\Rightarrow t^2-4t+3=0\Rightarrow$ \rest{1 s} (subiendo) y \rest{3 s} (bajando).\end{enumerate}


% =====================================================
\seccio{8}{Construcción de funciones a partir de un enunciado}

\exf{1} \begin{enumerate}[label=\alph*),leftmargin=*,itemsep=2pt,topsep=2pt]\item $\begin{array}{c|ccccc}x&0&1&2&3&4\\\hline y&5&7&9&11&13\end{array}$ \quad $\res{y=5+2x}$\item Una \rest{función lineal}: la gráfica es una recta que corta el eje $OY$ en 5 y sube 2 por cada hora.

\begin{center}\begin{tikzpicture}[x=0.376cm,y=0.16cm,baseline={(current bounding box.center)},font=\scriptsize]
\draw[muted!25,line width=.3pt] (0,-1)--(0,22);
\draw[muted!25,line width=.3pt] (1,-1)--(1,22);
\draw[muted!25,line width=.3pt] (2,-1)--(2,22);
\draw[muted!25,line width=.3pt] (3,-1)--(3,22);
\draw[muted!25,line width=.3pt] (4,-1)--(4,22);
\draw[muted!25,line width=.3pt] (5,-1)--(5,22);
\draw[muted!25,line width=.3pt] (6,-1)--(6,22);
\draw[muted!25,line width=.3pt] (7,-1)--(7,22);
\draw[muted!25,line width=.3pt] (8,-1)--(8,22);
\draw[muted!25,line width=.3pt] (-0.5,0)--(8,0);
\draw[muted!25,line width=.3pt] (-0.5,2)--(8,2);
\draw[muted!25,line width=.3pt] (-0.5,4)--(8,4);
\draw[muted!25,line width=.3pt] (-0.5,6)--(8,6);
\draw[muted!25,line width=.3pt] (-0.5,8)--(8,8);
\draw[muted!25,line width=.3pt] (-0.5,10)--(8,10);
\draw[muted!25,line width=.3pt] (-0.5,12)--(8,12);
\draw[muted!25,line width=.3pt] (-0.5,14)--(8,14);
\draw[muted!25,line width=.3pt] (-0.5,16)--(8,16);
\draw[muted!25,line width=.3pt] (-0.5,18)--(8,18);
\draw[muted!25,line width=.3pt] (-0.5,20)--(8,20);
\draw[muted!25,line width=.3pt] (-0.5,22)--(8,22);
\draw[-latex,muted] (-0.5,0)--(8,0);
\draw[-latex,muted] (0,-1)--(0,22);
\node[below=-1pt,text=muted] at (1,0) {$1$};
\node[below=-1pt,text=muted] at (2,0) {$2$};
\node[below=-1pt,text=muted] at (3,0) {$3$};
\node[below=-1pt,text=muted] at (4,0) {$4$};
\node[below=-1pt,text=muted] at (5,0) {$5$};
\node[below=-1pt,text=muted] at (6,0) {$6$};
\node[below=-1pt,text=muted] at (7,0) {$7$};
\node[below=-1pt,text=muted] at (8,0) {$8$};
\node[left=-1pt,text=muted] at (0,4) {$4$};
\node[left=-1pt,text=muted] at (0,8) {$8$};
\node[left=-1pt,text=muted] at (0,12) {$12$};
\node[left=-1pt,text=muted] at (0,16) {$16$};
\node[left=-1pt,text=muted] at (0,20) {$20$};
\begin{scope}\clip (-0.5,-1) rectangle (8,22);
\draw[accent,line width=1.1pt,line join=round] (0,5)--(0.064,5.128)--(0.128,5.256)--(0.192,5.384)--(0.256,5.512)--(0.32,5.64)--(0.384,5.768)--(0.448,5.896)--(0.512,6.024)--(0.576,6.152)--(0.64,6.28)--(0.704,6.408)--(0.768,6.536)--(0.832,6.664)--(0.896,6.792)--(0.96,6.92)--(1.024,7.048)--(1.088,7.176)--(1.152,7.304)--(1.216,7.432)--(1.28,7.56)--(1.344,7.688)--(1.408,7.816)--(1.472,7.944)--(1.536,8.072)--(1.6,8.2)--(1.664,8.328)--(1.728,8.456)--(1.792,8.584)--(1.856,8.712)--(1.92,8.84)--(1.984,8.968)--(2.048,9.096)--(2.112,9.224)--(2.176,9.352)--(2.24,9.48)--(2.304,9.608)--(2.368,9.736)--(2.432,9.864)--(2.496,9.992)--(2.56,10.12)--(2.624,10.248)--(2.688,10.376)--(2.752,10.504)--(2.816,10.632)--(2.88,10.76)--(2.944,10.888)--(3.008,11.016)--(3.072,11.144)--(3.136,11.272)--(3.2,11.4)--(3.264,11.528)--(3.328,11.656)--(3.392,11.784)--(3.456,11.912)--(3.52,12.04)--(3.584,12.168)--(3.648,12.296)--(3.712,12.424)--(3.776,12.552)--(3.84,12.68)--(3.904,12.808)--(3.968,12.936)--(4.032,13.064)--(4.096,13.192)--(4.16,13.32)--(4.224,13.448)--(4.288,13.576)--(4.352,13.704)--(4.416,13.832)--(4.48,13.96)--(4.544,14.088)--(4.608,14.216)--(4.672,14.344)--(4.736,14.472)--(4.8,14.6)--(4.864,14.728)--(4.928,14.856)--(4.992,14.984)--(5.056,15.112)--(5.12,15.24)--(5.184,15.368)--(5.248,15.496)--(5.312,15.624)--(5.376,15.752)--(5.44,15.88)--(5.504,16.008)--(5.568,16.136)--(5.632,16.264)--(5.696,16.392)--(5.76,16.52)--(5.824,16.648)--(5.888,16.776)--(5.952,16.904)--(6.016,17.032)--(6.08,17.16)--(6.144,17.288)--(6.208,17.416)--(6.272,17.544)--(6.336,17.672)--(6.4,17.8)--(6.464,17.928)--(6.528,18.056)--(6.592,18.184)--(6.656,18.312)--(6.72,18.44)--(6.784,18.568)--(6.848,18.696)--(6.912,18.824)--(6.976,18.952)--(7.04,19.08)--(7.104,19.208)--(7.168,19.336)--(7.232,19.464)--(7.296,19.592)--(7.36,19.72)--(7.424,19.848)--(7.488,19.976)--(7.552,20.104)--(7.616,20.232)--(7.68,20.36)--(7.744,20.488)--(7.808,20.616)--(7.872,20.744)--(7.936,20.872)--(8,21);
\filldraw[fill=blue!70!black,draw=blue!70!black,line width=.8pt] (0,5) circle (1.5pt);
\filldraw[fill=blue!70!black,draw=blue!70!black,line width=.8pt] (4,13) circle (1.5pt);
\filldraw[fill=blue!70!black,draw=blue!70!black,line width=.8pt] (7,19) circle (1.5pt);
\end{scope}\end{tikzpicture}\end{center}
\item $5+2x=19\Rightarrow\res{x=7}$ horas\end{enumerate}

\ex{1} \begin{enumerate}[label=\alph*),leftmargin=*,itemsep=2pt,topsep=2pt]\item $\begin{array}{c|cccccc}x&1&2&3&4&5&6\\\hline y&360&180&120&90&72&60\end{array}$ \quad $\res{y=\dfrac{360}{x}}$\item Solo \rest{números naturales} (1, 2, 3…): \rest{no se unen}. Es de proporcionalidad inversa.\item $360:7\approx\res{51{,}43\text{ €}}$ (y sobra algún céntimo).\end{enumerate}

\ex{2} \begin{enumerate}[label=\alph*),leftmargin=*,itemsep=2pt,topsep=2pt]\item $\res{y=1\,\,000+25x}$\item $1\,\,000+25x=6\,\,000\Rightarrow x=200$ min (\rest{3 h 20 min}). $D=\res{[0,\,200]}$, $R=\res{[1\,\,000;\ 6\,\,000]}$.\item $y(60)=\res{2\,\,500\text{ L}}$. \ $1\,\,000+25x=4\,\,000\Rightarrow x=\res{120}$ min (2 horas).\end{enumerate}

\ex{2} \begin{enumerate}[label=\alph*),leftmargin=*,itemsep=2pt,topsep=2pt]\item Gasta $6:100=0{,}06$ L por km: $\res{y=45-0{,}06x}$.\item $45-0{,}06\cdot300=\res{27\text{ L}}$\item Se acaba cuando $45-0{,}06x=0\Rightarrow x=750$ km: $D=\res{[0,\,750]}$, $R=\res{[0,\,45]}$.\end{enumerate}

\ex{2} $x\cdot y=24\Rightarrow\res{y=\dfrac{24}{x}}$ \quad $\begin{array}{c|cccccc}x&1&2&3&4&6&8\\\hline y&24&12&8&6&4&3\end{array}$ \ Es de \rest{proporcionalidad inversa} (si la base se duplica, la altura se reduce a la mitad). $D=\res{(0,\,+\infty)}$: la base debe ser positiva.

\begin{center}\begin{tikzpicture}[x=0.243cm,y=0.145cm,baseline={(current bounding box.center)},font=\scriptsize]
\draw[muted!25,line width=.3pt] (-1,-1)--(-1,26);
\draw[muted!25,line width=.3pt] (0,-1)--(0,26);
\draw[muted!25,line width=.3pt] (1,-1)--(1,26);
\draw[muted!25,line width=.3pt] (2,-1)--(2,26);
\draw[muted!25,line width=.3pt] (3,-1)--(3,26);
\draw[muted!25,line width=.3pt] (4,-1)--(4,26);
\draw[muted!25,line width=.3pt] (5,-1)--(5,26);
\draw[muted!25,line width=.3pt] (6,-1)--(6,26);
\draw[muted!25,line width=.3pt] (7,-1)--(7,26);
\draw[muted!25,line width=.3pt] (8,-1)--(8,26);
\draw[muted!25,line width=.3pt] (9,-1)--(9,26);
\draw[muted!25,line width=.3pt] (10,-1)--(10,26);
\draw[muted!25,line width=.3pt] (11,-1)--(11,26);
\draw[muted!25,line width=.3pt] (12,-1)--(12,26);
\draw[muted!25,line width=.3pt] (13,-1)--(13,26);
\draw[muted!25,line width=.3pt] (-1,0)--(13,0);
\draw[muted!25,line width=.3pt] (-1,2)--(13,2);
\draw[muted!25,line width=.3pt] (-1,4)--(13,4);
\draw[muted!25,line width=.3pt] (-1,6)--(13,6);
\draw[muted!25,line width=.3pt] (-1,8)--(13,8);
\draw[muted!25,line width=.3pt] (-1,10)--(13,10);
\draw[muted!25,line width=.3pt] (-1,12)--(13,12);
\draw[muted!25,line width=.3pt] (-1,14)--(13,14);
\draw[muted!25,line width=.3pt] (-1,16)--(13,16);
\draw[muted!25,line width=.3pt] (-1,18)--(13,18);
\draw[muted!25,line width=.3pt] (-1,20)--(13,20);
\draw[muted!25,line width=.3pt] (-1,22)--(13,22);
\draw[muted!25,line width=.3pt] (-1,24)--(13,24);
\draw[muted!25,line width=.3pt] (-1,26)--(13,26);
\draw[-latex,muted] (-1,0)--(13,0);
\draw[-latex,muted] (0,-1)--(0,26);
\node[below=-1pt,text=muted] at (2,0) {$2$};
\node[below=-1pt,text=muted] at (4,0) {$4$};
\node[below=-1pt,text=muted] at (6,0) {$6$};
\node[below=-1pt,text=muted] at (8,0) {$8$};
\node[below=-1pt,text=muted] at (10,0) {$10$};
\node[below=-1pt,text=muted] at (12,0) {$12$};
\node[left=-1pt,text=muted] at (0,4) {$4$};
\node[left=-1pt,text=muted] at (0,8) {$8$};
\node[left=-1pt,text=muted] at (0,12) {$12$};
\node[left=-1pt,text=muted] at (0,16) {$16$};
\node[left=-1pt,text=muted] at (0,20) {$20$};
\node[left=-1pt,text=muted] at (0,24) {$24$};
\begin{scope}\clip (-1,-1) rectangle (13,26);
\draw[accent,line width=1.1pt,line join=round] (0.924,25.968)--(1.021,23.506)--(1.118,21.471)--(1.215,19.76)--(1.311,18.301)--(1.408,17.043)--(1.505,15.947)--(1.602,14.983)--(1.699,14.129)--(1.795,13.367)--(1.892,12.684)--(1.989,12.066)--(2.086,11.506)--(2.183,10.996)--(2.279,10.529)--(2.376,10.1)--(2.473,9.705)--(2.57,9.339)--(2.667,9)--(2.763,8.685)--(2.86,8.391)--(2.957,8.116)--(3.054,7.859)--(3.151,7.618)--(3.247,7.391)--(3.344,7.177)--(3.441,6.975)--(3.538,6.784)--(3.635,6.603)--(3.731,6.432)--(3.828,6.269)--(3.925,6.115)--(4.022,5.967)--(4.119,5.827)--(4.215,5.693)--(4.312,5.566)--(4.409,5.443)--(4.506,5.326)--(4.603,5.214)--(4.699,5.107)--(4.796,5.004)--(4.893,4.905)--(4.99,4.81)--(5.087,4.718)--(5.183,4.63)--(5.28,4.545)--(5.377,4.463)--(5.474,4.385)--(5.571,4.308)--(5.667,4.235)--(5.764,4.164)--(5.861,4.095)--(5.958,4.028)--(6.055,3.964)--(6.151,3.902)--(6.248,3.841)--(6.345,3.783)--(6.442,3.726)--(6.539,3.671)--(6.635,3.617)--(6.732,3.565)--(6.829,3.514)--(6.926,3.465)--(7.023,3.418)--(7.119,3.371)--(7.216,3.326)--(7.313,3.282)--(7.41,3.239)--(7.507,3.197)--(7.603,3.156)--(7.7,3.117)--(7.797,3.078)--(7.894,3.04)--(7.991,3.004)--(8.087,2.968)--(8.184,2.932)--(8.281,2.898)--(8.378,2.865)--(8.475,2.832)--(8.571,2.8)--(8.668,2.769)--(8.765,2.738)--(8.862,2.708)--(8.959,2.679)--(9.055,2.65)--(9.152,2.622)--(9.249,2.595)--(9.346,2.568)--(9.443,2.542)--(9.539,2.516)--(9.636,2.491)--(9.733,2.466)--(9.83,2.442)--(9.927,2.418)--(10.023,2.394)--(10.12,2.371)--(10.217,2.349)--(10.314,2.327)--(10.411,2.305)--(10.507,2.284)--(10.604,2.263)--(10.701,2.243)--(10.798,2.223)--(10.895,2.203)--(10.991,2.184)--(11.088,2.164)--(11.185,2.146)--(11.282,2.127)--(11.379,2.109)--(11.475,2.091)--(11.572,2.074)--(11.669,2.057)--(11.766,2.04)--(11.863,2.023)--(11.959,2.007)--(12.056,1.991)--(12.153,1.975)--(12.25,1.959)--(12.347,1.944)--(12.443,1.929)--(12.54,1.914)--(12.637,1.899)--(12.734,1.885)--(12.831,1.871)--(12.927,1.857)--(13,1.846);
\filldraw[fill=blue!70!black,draw=blue!70!black,line width=.8pt] (2,12) circle (1.5pt);
\filldraw[fill=blue!70!black,draw=blue!70!black,line width=.8pt] (4,6) circle (1.5pt);
\filldraw[fill=blue!70!black,draw=blue!70!black,line width=.8pt] (8,3) circle (1.5pt);
\end{scope}\end{tikzpicture}\end{center}


\ex{2} \begin{enumerate}[label=\alph*),leftmargin=*,itemsep=2pt,topsep=2pt]\item El otro lado mide $\frac{60-2x}{2}=30-x$: $\res{A(x)=x(30-x)=30x-x^2}$.\item $D=\res{(0,\,30)}$. \ $A(5)=125$, $A(10)=200$, $A(15)=225$, $A(20)=200$, $A(25)=125$ (m$^2$).\item La gráfica es una parábola con el máximo en $(15,\,225)$: un \rest{cuadrado de 15 m de lado}, con 225 m$^2$.\end{enumerate}

\ex{2} \begin{enumerate}[label=\alph*),leftmargin=*,itemsep=2pt,topsep=2pt]\item $\res{y_A=10+0{,}05x}$, \ $\res{y_B=0{,}1x}$\item 100 min: A \rest{15 €}, B \rest{10 €}. \ 300 min: A \rest{25 €}, B \rest{30 €}.\item $10+0{,}05x=0{,}1x\Rightarrow x=\res{200}$ minutos (20 € las dos). Para menos de 200 minutos conviene la \rest{B}; para más, la \rest{A}.

\begin{center}\begin{tikzpicture}[x=0.009cm,y=0.052cm,baseline={(current bounding box.center)},font=\scriptsize]
\draw[muted!25,line width=.3pt] (0,-3)--(0,45);
\draw[muted!25,line width=.3pt] (50,-3)--(50,45);
\draw[muted!25,line width=.3pt] (100,-3)--(100,45);
\draw[muted!25,line width=.3pt] (150,-3)--(150,45);
\draw[muted!25,line width=.3pt] (200,-3)--(200,45);
\draw[muted!25,line width=.3pt] (250,-3)--(250,45);
\draw[muted!25,line width=.3pt] (300,-3)--(300,45);
\draw[muted!25,line width=.3pt] (350,-3)--(350,45);
\draw[muted!25,line width=.3pt] (400,-3)--(400,45);
\draw[muted!25,line width=.3pt] (-20,0)--(420,0);
\draw[muted!25,line width=.3pt] (-20,5)--(420,5);
\draw[muted!25,line width=.3pt] (-20,10)--(420,10);
\draw[muted!25,line width=.3pt] (-20,15)--(420,15);
\draw[muted!25,line width=.3pt] (-20,20)--(420,20);
\draw[muted!25,line width=.3pt] (-20,25)--(420,25);
\draw[muted!25,line width=.3pt] (-20,30)--(420,30);
\draw[muted!25,line width=.3pt] (-20,35)--(420,35);
\draw[muted!25,line width=.3pt] (-20,40)--(420,40);
\draw[muted!25,line width=.3pt] (-20,45)--(420,45);
\draw[-latex,muted] (-20,0)--(420,0);
\draw[-latex,muted] (0,-3)--(0,45);
\node[below=-1pt,text=muted] at (100,0) {$100$};
\node[below=-1pt,text=muted] at (200,0) {$200$};
\node[below=-1pt,text=muted] at (300,0) {$300$};
\node[below=-1pt,text=muted] at (400,0) {$400$};
\node[left=-1pt,text=muted] at (0,10) {$10$};
\node[left=-1pt,text=muted] at (0,20) {$20$};
\node[left=-1pt,text=muted] at (0,30) {$30$};
\node[left=-1pt,text=muted] at (0,40) {$40$};
\begin{scope}\clip (-20,-3) rectangle (420,45);
\draw[blue!70!black,line width=1.1pt,line join=round] (0,10)--(3.36,10.168)--(6.72,10.336)--(10.08,10.504)--(13.44,10.672)--(16.8,10.84)--(20.16,11.008)--(23.52,11.176)--(26.88,11.344)--(30.24,11.512)--(33.6,11.68)--(36.96,11.848)--(40.32,12.016)--(43.68,12.184)--(47.04,12.352)--(50.4,12.52)--(53.76,12.688)--(57.12,12.856)--(60.48,13.024)--(63.84,13.192)--(67.2,13.36)--(70.56,13.528)--(73.92,13.696)--(77.28,13.864)--(80.64,14.032)--(84,14.2)--(87.36,14.368)--(90.72,14.536)--(94.08,14.704)--(97.44,14.872)--(100.8,15.04)--(104.16,15.208)--(107.52,15.376)--(110.88,15.544)--(114.24,15.712)--(117.6,15.88)--(120.96,16.048)--(124.32,16.216)--(127.68,16.384)--(131.04,16.552)--(134.4,16.72)--(137.76,16.888)--(141.12,17.056)--(144.48,17.224)--(147.84,17.392)--(151.2,17.56)--(154.56,17.728)--(157.92,17.896)--(161.28,18.064)--(164.64,18.232)--(168,18.4)--(171.36,18.568)--(174.72,18.736)--(178.08,18.904)--(181.44,19.072)--(184.8,19.24)--(188.16,19.408)--(191.52,19.576)--(194.88,19.744)--(198.24,19.912)--(201.6,20.08)--(204.96,20.248)--(208.32,20.416)--(211.68,20.584)--(215.04,20.752)--(218.4,20.92)--(221.76,21.088)--(225.12,21.256)--(228.48,21.424)--(231.84,21.592)--(235.2,21.76)--(238.56,21.928)--(241.92,22.096)--(245.28,22.264)--(248.64,22.432)--(252,22.6)--(255.36,22.768)--(258.72,22.936)--(262.08,23.104)--(265.44,23.272)--(268.8,23.44)--(272.16,23.608)--(275.52,23.776)--(278.88,23.944)--(282.24,24.112)--(285.6,24.28)--(288.96,24.448)--(292.32,24.616)--(295.68,24.784)--(299.04,24.952)--(302.4,25.12)--(305.76,25.288)--(309.12,25.456)--(312.48,25.624)--(315.84,25.792)--(319.2,25.96)--(322.56,26.128)--(325.92,26.296)--(329.28,26.464)--(332.64,26.632)--(336,26.8)--(339.36,26.968)--(342.72,27.136)--(346.08,27.304)--(349.44,27.472)--(352.8,27.64)--(356.16,27.808)--(359.52,27.976)--(362.88,28.144)--(366.24,28.312)--(369.6,28.48)--(372.96,28.648)--(376.32,28.816)--(379.68,28.984)--(383.04,29.152)--(386.4,29.32)--(389.76,29.488)--(393.12,29.656)--(396.48,29.824)--(399.84,29.992)--(403.2,30.16)--(406.56,30.328)--(409.92,30.496)--(413.28,30.664)--(416.64,30.832)--(420,31);
\draw[orange!85!black,line width=1.1pt,line join=round] (0,0)--(3.36,0.336)--(6.72,0.672)--(10.08,1.008)--(13.44,1.344)--(16.8,1.68)--(20.16,2.016)--(23.52,2.352)--(26.88,2.688)--(30.24,3.024)--(33.6,3.36)--(36.96,3.696)--(40.32,4.032)--(43.68,4.368)--(47.04,4.704)--(50.4,5.04)--(53.76,5.376)--(57.12,5.712)--(60.48,6.048)--(63.84,6.384)--(67.2,6.72)--(70.56,7.056)--(73.92,7.392)--(77.28,7.728)--(80.64,8.064)--(84,8.4)--(87.36,8.736)--(90.72,9.072)--(94.08,9.408)--(97.44,9.744)--(100.8,10.08)--(104.16,10.416)--(107.52,10.752)--(110.88,11.088)--(114.24,11.424)--(117.6,11.76)--(120.96,12.096)--(124.32,12.432)--(127.68,12.768)--(131.04,13.104)--(134.4,13.44)--(137.76,13.776)--(141.12,14.112)--(144.48,14.448)--(147.84,14.784)--(151.2,15.12)--(154.56,15.456)--(157.92,15.792)--(161.28,16.128)--(164.64,16.464)--(168,16.8)--(171.36,17.136)--(174.72,17.472)--(178.08,17.808)--(181.44,18.144)--(184.8,18.48)--(188.16,18.816)--(191.52,19.152)--(194.88,19.488)--(198.24,19.824)--(201.6,20.16)--(204.96,20.496)--(208.32,20.832)--(211.68,21.168)--(215.04,21.504)--(218.4,21.84)--(221.76,22.176)--(225.12,22.512)--(228.48,22.848)--(231.84,23.184)--(235.2,23.52)--(238.56,23.856)--(241.92,24.192)--(245.28,24.528)--(248.64,24.864)--(252,25.2)--(255.36,25.536)--(258.72,25.872)--(262.08,26.208)--(265.44,26.544)--(268.8,26.88)--(272.16,27.216)--(275.52,27.552)--(278.88,27.888)--(282.24,28.224)--(285.6,28.56)--(288.96,28.896)--(292.32,29.232)--(295.68,29.568)--(299.04,29.904)--(302.4,30.24)--(305.76,30.576)--(309.12,30.912)--(312.48,31.248)--(315.84,31.584)--(319.2,31.92)--(322.56,32.256)--(325.92,32.592)--(329.28,32.928)--(332.64,33.264)--(336,33.6)--(339.36,33.936)--(342.72,34.272)--(346.08,34.608)--(349.44,34.944)--(352.8,35.28)--(356.16,35.616)--(359.52,35.952)--(362.88,36.288)--(366.24,36.624)--(369.6,36.96)--(372.96,37.296)--(376.32,37.632)--(379.68,37.968)--(383.04,38.304)--(386.4,38.64)--(389.76,38.976)--(393.12,39.312)--(396.48,39.648)--(399.84,39.984)--(403.2,40.32)--(406.56,40.656)--(409.92,40.992)--(413.28,41.328)--(416.64,41.664)--(420,42);
\filldraw[fill=violet,draw=violet,line width=.8pt] (200,20) circle (1.5pt);
\node[text=blue!70!black] at (400,26) {A};
\node[text=orange!85!black] at (330,37) {B};
\end{scope}\end{tikzpicture}\end{center}
\end{enumerate}

\ex{3} El 5\,\% de $x$ es $0{,}05x$: $\res{y=900+0{,}05x}$. \ $y(6\,\,000)=900+300=\res{1\,\,200\text{ €}}$. \ $900+0{,}05x=1\,\,500\Rightarrow x=\dfrac{600}{0{,}05}=\res{12\,\,000\text{ €}}$.

\ex{2} \begin{enumerate}[label=\alph*),leftmargin=*,itemsep=2pt,topsep=2pt]\item Cada entrada vale \rest{5 €}; el gasto fijo es $8-5=$ \rest{3 €}.\item $\res{y=5x+3}$; \ $y(10)=\res{53\text{ €}}$\item $5x+3=98\Rightarrow\res{x=19}$ entradas. \rest{No}: el número de entradas es un número natural.\end{enumerate}


\end{document}
